% !TeX TXS-program:bibliography = txs:///bibtex

\documentclass{amsart}

\usepackage{amsfonts, amsmath, amssymb, pifont}
\usepackage{mathtools}
\usepackage{hyperref}
\usepackage{enumitem}
\usepackage{todonotes}

\usepackage{xcolor}
\usepackage{tikz}
\usetikzlibrary{calc}

\newlength{\superthick}
\newlength{\cornerradius}
%\setlength{\superthick}{2.4pt}
\setlength{\superthick}{0.08cm}
\setlength{\cornerradius}{5pt}
\tikzstyle{corner}=[rounded corners=\cornerradius]
\tikzstyle{string}=[line width=\superthick]
%\tikzstyle{dot}=[circle, inner sep=0pt, minimum size=4.8pt]
\tikzstyle{dot}=[circle, inner sep=0pt, minimum size=7pt]

\theoremstyle{plain}
\newtheorem{thm}{Theorem}[section]
\newtheorem*{thm*}{Theorem}
\newtheorem{clm}[thm]{Claim}
\newtheorem{lem}[thm]{Lemma}
\newtheorem*{lem*}{Lemma}
\newtheorem{conj}[thm]{Conjecture}
\newtheorem*{conj*}{Conjecture}
\newtheorem{cor}[thm]{Corollary}
\newtheorem*{cor*}{Corollary}
\newtheorem{prop}[thm]{Proposition}
\newtheorem{thmletter}{Theorem}[section]
\renewcommand*{\thethmletter}{\Alph{thmletter}}


\theoremstyle{definition}
\newtheorem{defn}[thm]{Definition}
\newtheorem{exam}[thm]{Example}
\newtheorem{notn}[thm]{Notation}
\newtheorem{constr}[thm]{Construction}
%
\theoremstyle{remark}
\newtheorem{rem}[thm]{Remark}

\numberwithin{equation}{section}

\newcommand{\defnemph}[1]{\emph{#1}}

\newcommand{\NN}{\mathbb{N}}
\newcommand{\ZZ}{\mathbb{Z}}
\newcommand{\QQ}{\mathbb{Q}}
\newcommand{\RR}{\mathbb{R}}
\newcommand{\CC}{\mathbb{C}}
\newcommand{\KK}{\mathbb{K}}
\newcommand{\ident}{{\rm id}}

\DeclareMathOperator{\Hom}{Hom}
\DeclareMathOperator{\End}{End}
\DeclareMathOperator{\Sym}{Sym}


\begin{document}

\title[Remarks on realizations]{Some remarks on realizations defining diagrammatic Hecke categories}
\author{Amit Hazi}

\email{A.Hazi@leeds.ac.uk}

\begin{abstract}
Every diagrammatic Hecke category is constructed from a realization, which generalizes the notion of a reflection representation of a Coxeter group. 
Common assumptions on realizations to ensure that the resulting categories are well behaved include Demazure surjectivity and the parabolic property (for anti-spherical quotients). 
We show that these assumptions are in fact unnecessary. 
%We show that the analogues of Soergel's categorification theorem always hold for the diagrammatic Hecke category and the diagrammatic anti-spherical category, even in the absence of Demazure surjectivity and the parabolic property. 
We also give a non-inductive description of the rotational scalars for unbalanced realizations, as well as a simpler criterion to check for balancedness. 
\end{abstract}

\subjclass[2020]{20C08 (primary); 20G15, 20F55 (secondary).}

\maketitle

\section*{Introduction}

Let $(W,S)$ be a Coxeter system, and let $\mathbb{H}:=\mathbb{H}(W,S)$ be the corresponding Iwahori--Hecke algebra. 
A \defnemph{Hecke category} is a graded additive monoidal category $\mathcal{H}$ which categorifies $\mathbb{H}$. 
Such categories (and their categorical actions) are ubiquitous in Lie theory. 
In characteristic $0$ their behavior is often determined by Kazhdan--Lusztig theory, which gives insight into how Kazhdan--Lusztig polynomials control the categories they act on (e.g.~the Kazhdan--Lusztig conjectures for the BGG category $\mathcal{O}$).
In positive characteristic, Hecke categories are used to \emph{define} $p$-Kazhdan--Lusztig theory, providing a replacement for Kazhdan--Lusztig polynomials with applications to modular representation theory. 

%Hecke categories appear in various guises throughout (geometric) representation theory, but the most general constructions are diagrammatic, i.e.~essentially a presentation of a category by generators and relations. 
As with the Hecke algebra itself, the most general construction of a Hecke category is via a (diagrammatic) presentation in terms of generators and relations due to Elias--Williamson \cite{EW-SoergelCalc}. 
The fundamental input data for such a presentation is a \defnemph{realization} of the Coxeter system $(W,S)$. 
Realizations generalize the notion of a reflection representation (specifically that of a \defnemph{reflection faithful} representation) with labeled roots and coroots. 
Such representations are the foundation for Soergel's original algebraic categorification of the Hecke algebra \cite{soergel-unzerlegbareBimoduln}.
In the diagrammatic setting, realizations were first introduced by Elias in \cite{dihedralcathedral}. 

In many applications the underlying realization is the action of a (finite or affine) Weyl group on a Cartan subalgebra, which is almost always well behaved. % in characteristic $0$.
However, finding the correct completely general definition is surprisingly subtle. 
For example, 7 years after introducing general realizations Elias--Williamson showed that the original definition contained an incorrect inference \cite{ew-loccalc}; this error was subsequently corrected by the author in \cite{2colJW}. 
Moreover, in the literature the relatively weak condition of \defnemph{Demazure surjectivity} is often freely assumed to ensure that the corresponding diagrammatic category actually categorifies the Hecke algebra. 
%Neither of these assumptions is particularly strong or difficult to explain. 
For \defnemph{anti-spherical} quotients of Hecke categories, Libedinsky--Williamson require an additional strong assumption (the \defnemph{parabolic property}) which often fails in positive characteristic \cite{LW-antispher}. 
Both of these conditions add complexity to an already complicated theory.

The goal of this paper is to simplify the theory of diagrammatic Hecke categories by showing that these assumptions are broadly unnecessary. 
Our main results are as follows.

\begin{thmletter}[{Corollary \ref{cor:cinchedSoergelcatthm}}] \label{thm:intro1}
The diagrammatic category constructed by Elias--Williamson always admits a quotient which categorifies the Hecke algebra.
\end{thmletter}

\begin{thmletter}[{Corollary \ref{cor:antispherSoergelcatthm}}] \label{thm:intro2}
The anti-spherical quotient of a Hecke category (\textit{\`{a} la} Libedinsky--Williamson) always categorifies the anti-spherical module for the Hecke algebra.
\end{thmletter}

In addition, we provide an easier way to check for \defnemph{balancedness}, another assumption on the realization which simplifies some of the notation in the diagrammatic Hecke category.

\begin{thmletter}[{Corollary \ref{cor:balancedcriterion}}]
A realization is balanced if and only if for all distinct $s,t \in S$ for which $st$ has odd order $m$, the $(s,t)$th entry of the Cartan matrix is a root of the minimal polynomial of $-2\cos(\pi/m)$ over $\ZZ$. 
\end{thmletter}

This follows from a calculation of the general rotational scalar $[m-1]_s$ in Theorem \ref{thm:mminus1twocol}, which has applications to unbalanced realizations.

%Theorems \ref{thm:intro1} and \ref{thm:intro2} have antecedents in the literature, and in most cases
%
%
%In this paper we seek to correct some of these misconceptions, which . 
%Some slogans to guide us:
%\begin{itemize}
%\item Realizations are (almost) Cartan matrices. 
%
%More precisely, one can check that a collection of vectors and covectors forms a realization just by checking their pairing (Proposition \ref{prop:rlz=cartanmat}), and many properties of interest (e.g.~canonical bases) depend only on the Cartan matrix.
%
%\item The parabolic property \cite[\S 2.3]{LW-antispher} is unnecessary. 
%
%More precisely, \emph{any} anti-spherical quotient of a Hecke category gives a categorification of the corresponding anti-spherical module (Theorem \ref{thm:antispherLLbasis}).
%
%\item Demazure surjectivity is unnecessary. 
%
%More precisely, every Elias--Williamson diagrammatic category admits a (possibly trivial) quotient which categorifies the Hecke algebra (Theorem \ref{thm:LLcinched}).
%
%\item $2$-colored quantum numbers are unnecessary. 
%
%More precisely, we give a direct (i.e.~non-inductive) definition of the $2$-colored scalars $[m-1]_s$ and $[m-1]_t$ for odd $m=m_{st}$ in the unbalanced case (Theorem \ref{thm:mminus1twocol}). 
%
%\item Localization is a presentation of the Hecke category. 
%
%More precisely, the Hecke category can always be presented
%\footnote{Strictly speaking, we do not always have a matrix \emph{representation}; see below for details.}
%using matrices over the fraction field of \emph{some} polynomial ring (Remark \ref{rem:locpres}). 
%
%
%\end{itemize}

\subsection*{Acknowledgments}

The author would like to thank Chris Bowman, Maud De Visscher, and Emily Norton for inspiring this paper through many stimulating discussions on similar topics while working on \cite{BHN-modularWeylKac} and \cite{bdhn-hsp}.
The author is grateful for financial support from EPSRC (EP/W007509/1).

\section{Realizations from Cartan matrices}


In this section we will give a simplification of the usual definition of a realization, which will help us prove many of our later results in this paper. 

\begin{defn} \label{defn:rlz}
Let $\Bbbk$ be a commutative integral domain, and let $(W,S)$ be a Coxeter system.
\begin{enumerate}
 \item A \defnemph{$\Bbbk$-valued Cartan matrix} for $(W,S)$ is an $S \times S$ matrix $(a_{st})_{s,t \in S}$ with entries in $\Bbbk$ satisfying the following conditions:
 \begin{enumerate}[label=(\roman*)]
 \item $a_{ss}=2$ for all $s \in S$;
 
 \item For all distinct $s,t \in S$ such that $m_{st}<\infty$,
 \begin{equation} \label{eq:minpolycond}
 \begin{aligned}
 \Psi_{m_{st}}(a_{st} a_{ts})& =0 & &\text{if $m_{st}>2$,} \\
 a_{st}=a_{ts}& =0 & & \text{if $m_{st}=2$,}
 \end{aligned}
 \end{equation}
 where $\Psi_{m_{st}}$ denotes the minimal polynomial of $4\cos^2(\pi/m_{st})$ over $\QQ$ (see \cite[(15)]{2colJW}).
% \item $\Psi_{m_{st}}(a_{st}a_{ts})=0$ for all distinct $s,t \in S$ with $2<m_{st}<\infty$, where $\Psi_{m_{st}}$ denotes the minimal polynomial of $4\cos^2(\pi/m_{st})$ over $\QQ$ (see \cite[(15)]{2colJW}).
 \end{enumerate}
 
 \item A \defnemph{realization} of $(W,S)$ over $\Bbbk$ consists of a free finite-rank $\Bbbk$-module $V$ along with subsets $\{\alpha_s : s \in S\} \subset V$ and $\{\alpha_s^\vee : s \in S\} \subset V^* = \Hom(V,\Bbbk)$ (called simple roots and coroots respectively) such that the matrix $(\langle \alpha_s^\vee, \alpha_t\rangle)_{s,t \in S}$ is a $\Bbbk$-valued Cartan matrix of $(W,S)$.
 
%\item Suppose $(V,\{\alpha_s\},\{\alpha_s^\vee\})$ is a realization. 
%Let $U=\sum_s \Bbbk \alpha_s$ denote the span of the simple roots, and identify its dual $U^\ast$ with the quotient $V^\ast / U^\circ$ (where $U^\circ \leq V^\ast$ denotes the annihilator of $U$).
%The \defnemph{root subrealization} is the realization defined by $(U,\{\alpha_s\},\{\alpha_s^\vee+U^\circ\})$.
 
 \item Let $(V,\{\alpha_s\},\{\alpha_s^\vee\})$ and $(V',\{\alpha'_s\},\{(\alpha'_s)^\vee\})$ be realizations of $(W,S)$ over $\Bbbk$. 
 A \defnemph{morphism of realizations} is a $\Bbbk$-linear map $\phi:V \rightarrow V'$ such that $\phi(\alpha_s)=\alpha'_s$ and $\phi^\ast((\alpha'_s)^\vee)=\alpha_s^\vee$ for all $s \in S$. 
 Let $U \leq V'$ denote the image of $\phi$, and identify its dual $U^\ast$ with the quotient $V'^\ast/U^\circ$ (where $U^\circ \leq V'^\ast$ denotes the annihilator of $U$). 
 The \defnemph{image realization} is given by $(U,\{\alpha'_s\},\{(\alpha'_s)^\vee+U^\circ\})$. 
 Dually, let $K \leq V$ denote the kernel of $\phi$. 
 The \defnemph{coimage realization} is given by $(V/K,\{\alpha_s+K\},\{\alpha_s^\vee\})$, and is isomorphic to the image realization.
\end{enumerate}
\end{defn}

\begin{defn} \label{defn:2colqnum}
Let $A=\ZZ[x_s,x_t]$ be the integral polynomial ring in two variables. 
The \defnemph{(generic) two-colored quantum numbers} are defined as follows. 
First set $[0]_s=[0]_t=0$, $[1]_{s}=[1]_{t}=1$, $[2]_{s}=x_s$, and $[2]_{t}=x_t$. 
For $n>1$ we inductively define
\begin{align}
[n+1]_{s}& =[2]_{s} [n]_{t} - [n-1]_{s} \text{,} & [n+1]_{t} & =[2]_{t} [n]_{s} - [n-1]_{t} \text{.} \label{eq:twocolqnum}
\end{align}

Let $(a_{st})_{s,t \in S}$ be a $\Bbbk$-valued Cartan matrix for a Coxeter system $(W,S)$. 
For distinct $s,t \in S$ we define the corresponding two-colored quantum numbers in $\Bbbk$ by specializing to $x_s=-a_{st}$ and $x_t=-a_{ts}$.
\end{defn}

\begin{rem} \hfill
\begin{enumerate}
\item For any positive integer $m$, the real numbers $2\cos(\pi/m)$ and $4\cos^2(\pi/m)$ are algebraic integers, so $\Psi_m$ in fact has integer coefficients, and thus makes sense over any ring $\Bbbk$. 

\item Equation \eqref{eq:minpolycond} is equivalent to Abe's condition \cite[Assumption~1.1]{abe-homBS} on the vanishing of certain two-colored quantum binomial coefficients \cite[Theorem~3.8]{2colJW}. 
We feel that the minimal polynomial vanishing condition is easier to state and check in practice, especially for crystallographic Coxeter groups. 
%
%\item Realizations carry a natural notion of duality exchanging roots and coroots. 
%Our definitions are all self-dual. 
%In particular, the dual notion of the image of a morphism of realizations is the coimage, which is isomorphic.
\end{enumerate}
\end{rem}

In Definition \ref{defn:rlz} we have removed the usual condition that a realization should be a representation of $W$. 
Somewhat surprisingly this turns out to be unnecessary. 
This observation is not original but appears on the whole to have been overlooked.
For completeness we will state and prove it carefully.

\begin{prop}[{cf.~\cite[\S A.2]{dihedralcathedral}}] \label{prop:rlz=cartanmat}
Let $V$ be a realization as defined above. 
Then the assignments
\begin{align*}
s(x)& =x-\langle \alpha_s^\vee, x\rangle \alpha_s \text{,} & s(y)& =y-\langle y,\alpha_s \rangle \alpha_s^\vee
\end{align*}
for any $s \in S$, $x \in V$, and $y \in V^\ast$ define dual representations of $W$ on $V$ and $V^\ast$ respectively.
\end{prop}

\begin{proof}
We will first show that the first assignment defines a representation of $W$ on $V$. 
It is enough to check the defining relations of the Coxeter group on the action. 

Clearly for all $s \in S$ and $x \in V$ we have
\begin{align*}
s(s(x))& =s(x-\langle \alpha_s^\vee, x\rangle \alpha_s) \\
& =(x-\langle \alpha_s^\vee, x\rangle \alpha_s)-\langle \alpha_s^\vee, x\rangle (\alpha_s-\langle \alpha_s^\vee,\alpha_s\rangle \alpha_s) \\
& =(x-\langle \alpha_s^\vee, x\rangle \alpha_s)-\langle \alpha_s^\vee, x\rangle (\alpha_s-2\alpha_s) \\
& =(x-\langle \alpha_s^\vee, x\rangle \alpha_s)+\langle \alpha_s^\vee, x\rangle \alpha_s \\
& =x \text{.}
\end{align*}

%For distinct $s,t \in S$ let us define \defnemph{two-colored quantum numbers} as follows.
%First set $[0]_s=[0]_t=0$, $[1]_{s}=[1]_{t}=1$, $[2]_{s}=-\langle \alpha_s^\vee, \alpha_t\rangle$, and $[2]_{t}=-\langle \alpha_t^\vee, \alpha_s \rangle$. 
%For $n>1$ we inductively define
%\begin{align}
%[n+1]_{s}& =[2]_{s} [n]_{t} - [n-1]_{s} \text{,} & [n+1]_{t} & =[2]_{t} [n]_{s} - [n-1]_{t} \text{.} \label{eq:twocolqnum}
%\end{align}

Now suppose $s,t \in S$ are distinct with $m_{st}<\infty$. 
Write $k_t$ for the unique word $\dotsm tst$ which is an alternating product of $s$ and $t$ of length $k$ which ends with $t$.
For any $x \in V$ and any integer $k \geq 0$, we claim that
\begin{align*}
(2k)_t(x) & =x - [k]_s [k]_t \langle \alpha_s^\vee, x\rangle \alpha_s - [k]_s [k+1]_s \langle \alpha_t^\vee, x\rangle \alpha_s \\
& \quad \quad {} - [k]_t [k-1]_t \langle \alpha_s^\vee, x\rangle \alpha_t - [k]_s [k]_t \langle \alpha_t^\vee, x\rangle \alpha_t \text{,} \\
(2k+1)_t(x)& =x - [k]_s [k]_t \langle \alpha_s^\vee, x\rangle \alpha_s - [k]_s [k+1]_s \langle \alpha_t^\vee, x\rangle \alpha_s \\
& \quad \quad {} - [k]_t [k+1]_t \langle \alpha_s^\vee, x\rangle \alpha_t - [k+1]_s [k+1]_t \langle \alpha_t^\vee, x\rangle \alpha_t \text{.}
\end{align*}
We will prove this by induction. 

The base cases with $k=0$ are easy to check.
Suppose the result holds for $2k$. 
Then we have
\begin{align*}
(2k+1)_t(x)& =t(2k)_t(x) \\
& =t(x - [k]_s [k]_t \langle \alpha_s^\vee, x\rangle \alpha_s - [k]_s [k+1]_s \langle \alpha_t^\vee, x\rangle \alpha_s \\
& \quad \quad {} - [k]_t [k-1]_t \langle \alpha_s^\vee, x\rangle \alpha_t - [k]_s [k]_t \langle \alpha_t^\vee, x\rangle \alpha_t) \\
& =x - [k]_s [k]_t \langle \alpha_s^\vee, x\rangle \alpha_s - [k]_s [k+1]_s \langle \alpha_t^\vee, x\rangle \alpha_s \\
& \quad \quad {} + ([k]_t [k-1]_t - [2]_t [k]_s [k]_t)\langle \alpha_s^\vee, x\rangle \alpha_t \\
& \quad \quad {} + ([k]_s [k]_t - [2]_t [k]_s [k+1]_s - 1)\langle \alpha_t^\vee, x\rangle \alpha_t \\
& =x - [k]_s [k]_t \langle \alpha_s^\vee, x\rangle \alpha_s - [k]_s [k+1]_s \langle \alpha_t^\vee, x\rangle \alpha_s \\
& \quad \quad {} - [k]_t [k+1]_t \langle \alpha_s^\vee, x\rangle \alpha_t - ([k]_s [k+2]_t + 1)\langle \alpha_t^\vee, x\rangle \alpha_t \\
& =x - [k]_s [k]_t \langle \alpha_s^\vee, x\rangle \alpha_s - [k]_s [k+1]_s \langle \alpha_t^\vee, x\rangle \alpha_s \\
& \quad \quad {} - [k]_t [k+1]_t \langle \alpha_s^\vee, x\rangle \alpha_t - [k+1]_s [k+1]_t \langle \alpha_t^\vee, x\rangle \alpha_t \text{.}
\end{align*}
Here we have used the definition of two-colored quantum numbers, as well as the identity
\begin{equation}
[k]_s [k+2]_t + 1 = [k+1]_s [k+1]_t \text{.}
\end{equation}

Similarly, if the result holds for $2k+1$, then we have
\begin{align*}
(2k+2)_t(x)& =s(2k+1)_t(x) \\
& =s(x - [k]_s [k]_t \langle \alpha_s^\vee, x\rangle \alpha_s - [k]_s [k+1]_s \langle \alpha_t^\vee, x\rangle \alpha_s \\
& \quad \quad {} - [k]_t [k+1]_t \langle \alpha_s^\vee, x\rangle \alpha_t - [k+1]_s [k+1]_t \langle \alpha_t^\vee, x\rangle \alpha_t) \\
& =x + ([k]_s [k]_t - [2]_s [k]_t [k+1]_t - 1)\langle \alpha_s^\vee, x\rangle \alpha_s \\
& \quad \quad {} + ([k]_s [k+1]_s - [2]_s [k+1]_s [k+1]_t)\langle \alpha_t^\vee, x\rangle \alpha_s \\
& \quad \quad {} - [k]_t [k+1]_t \langle \alpha_s^\vee, x\rangle \alpha_t - [k+1]_s [k+1]_t \langle \alpha_t^\vee, x\rangle \alpha_t \\
& =x - ([k+2]_s [k]_t + 1)\langle \alpha_s^\vee, x\rangle \alpha_s - [k+2]_s [k+1]_s \langle \alpha_t^\vee, x\rangle \alpha_s \\
& \quad \quad {} - [k]_t [k+1]_t \langle \alpha_s^\vee, x\rangle \alpha_t - [k+1]_s [k+1]_t \langle \alpha_t^\vee, x\rangle \alpha_t \\
& =x - [k+1]_s [k+1]_t \langle \alpha_s^\vee, x\rangle \alpha_s - [k+2]_s [k+1]_s \langle \alpha_t^\vee, x\rangle \alpha_s \\
& \quad \quad {} - [k]_t [k+1]_t \langle \alpha_s^\vee, x\rangle \alpha_t - [k+1]_s [k+1]_t \langle \alpha_t^\vee, x\rangle \alpha_t \text{.}
\end{align*}

Since $(st)^m = (2m)_t$, the vanishing of $[m]_s$ and $[m]_t$ (a consequence of \eqref{eq:minpolycond} by \cite[Lemma~3.2]{2colJW}) implies that $(st)^m$ acts trivially on $V$. 
Thus $W$ acts on $V$ as prescribed.

To complete the proof, we note that for any $s \in S$, $x \in V$, and $y \in V^\ast$ we have
\begin{equation*}
\langle y, s(x) \rangle = \langle y, x \rangle - \langle y,\alpha_s \rangle \langle \alpha_s^\vee, x\rangle = \langle s(y), x\rangle \text{,}
\end{equation*}
so $W$ acts on $V^\ast$ as prescribed, in a dual manner to the action on $V$.
\end{proof}

\begin{cor}
Realizations as defined in Definition \ref{defn:rlz} are realizations in the sense of \cite[Definition~5.1]{2colJW}. 
In particular, a morphism of realizations of $(W,S)$ over $\Bbbk$ is always a homomorphism of $W$-representations.
\end{cor}

%\begin{proof}
%If $\phi:V \rightarrow V'$ is a morphism of realizations, then for any $s \in S$ and $x \in V$ we check that
%\begin{equation*}
%\phi(s(x))=\phi(x-\langle \alpha_s^\vee,x \rangle \alpha_s)=\phi(x)-\langle \phi^\ast((\alpha'_s)^\vee),x \rangle \phi(\alpha_s)=\phi(x)-\langle (\alpha'_s)^\vee,\phi(x)\rangle \alpha'_s=s(\phi(x)) \text{.}
%\end{equation*}
%\end{proof}

The most important consequence of this result is the construction of realizations directly from Cartan matrices. 
%Another consequence is the well-definedness of the following realizations,

\begin{defn}
Let $A=(a_{st})_{s,t \in S}$ be a $\Bbbk$-valued Cartan matrix for a Coxeter system $(W,S)$.
\begin{itemize}
\item The \defnemph{adjoint realization} $(V_{A,{\rm ad}}, \{\alpha_s\}, \{\alpha_s^\vee\})$ of $(W, S)$ with respect to $A$ is defined as follows. 
Let $V_{A,{\rm ad}}$ be a free $\Bbbk$-module with basis $\{\alpha_s : s \in S\}$ and define $\{\alpha_s^\vee\} \subset V_{A,{\rm ad}}^\ast$ by
\begin{equation}
\langle \alpha_s^\vee, \alpha_t \rangle = a_{st} \qquad \text{ for all $s,t \in S$.} \label{eq:rootpairingcartan}
\end{equation}

\item The \defnemph{simply-connected realization} $(V_{A,{\rm sc}}, \{\alpha_s\}, \{\alpha_s^\vee\})$ of $(W, S)$ with respect to $A$ is defined as follows. 
Abusing notation let $V_{A,{\rm sc}}^\ast$ be a free $\Bbbk$-module with basis $\{\alpha_s^\vee : s \in S\}$ and define $V_{A,{\rm sc}}=(V_{A,{\rm sc}}^\ast)^\ast$ and $\{\alpha_s\} \subset V_{A,{\rm sc}}$ by \eqref{eq:rootpairingcartan} above.
\end{itemize}
\end{defn}

Clearly the adjoint and simply-connected realizations with respect to a given Cartan matrix are dual to each other. 
The adjoint realization was called the ``universal realization'' in \cite[Remark~1.5]{BHN-modularWeylKac}; the newer names more closely match similar terminology from the theory of root data.

\begin{lem} \label{lem:adscuniversality}
For any realization $V$ of $(W, S)$ over $\Bbbk$ with Cartan matrix $A=(a_{st})_{s,t \in S}$, there are unique morphisms of realizations 
\begin{equation*}
V_{A,{\rm ad}} \longrightarrow V \text{,} \qquad \qquad \qquad  V \longrightarrow V_{A,{\rm sc}} \text{.}
\end{equation*}
\end{lem}

We call the image of the first morphism the \defnemph{root subrealization} of $V$, and the coimage of the second morphism the \defnemph{coroot quotient realization}.

\begin{proof}
We prove the adjoint case only as the argument for the simply-connected case is dual.
To avoid notational confusion let $\{\beta_s\}$ and $\{\beta_s^\vee\}$ denote the roots and coroots of $V$ respectively.
By construction there is only one linear map $\phi: V_{A,{\rm ad}} \rightarrow V$ mapping roots to roots. 
Moreover for all $s,t \in S$ we have
\begin{equation*}
\langle \phi^\ast(\beta_s^\vee),\alpha_t \rangle=\langle \beta_s^\vee,\phi(\alpha_t)\rangle=\langle \beta_s^\vee,\beta_t\rangle=a_{st}=\langle \alpha_s^\vee, \alpha_t \rangle \text{.}
\end{equation*}
As $\alpha_s^\vee \in V_{A,{\rm ad}}^\ast$ is defined by this condition it follows that $\phi^\ast(\beta_s^\vee)=\alpha_s^\vee$.
\end{proof}


Realizations of $(W,S)$ over $\Bbbk$ naturally form a category, and later we will see that the construction of the Hecke category from a realization is functorial. 

\section{Unbalanced realizations without quantum numbers}

Recall that a realization of a Coxeter system $(W,S)$ is called \defnemph{balanced} if for all distinct $s,t \in S$ with $m_{st}<\infty$, the corresponding two-colored quantum numbers $[m_{st}-1]_s$ and $[m_{st}-1]_t$ are both $1$. 
The majority of the literature on diagrammatic Hecke categories solely treats balanced realizations, as they are notationally more straightforward to work with. 
To check for balancedness it is enough to check when $m_{st}$ is odd, as for even $m_{st}$ it can be shown that $[m_{st}-1]_s=[m_{st}-1]_t=1$ automatically \cite[Lemma~6.25]{ew-loccalc}.

In the general unbalanced case, the scalars $[m_{st}-1]_s$ and $[m_{st}-1]_t$ play a role in the diagrammatic relations defining the Hecke category. 
In this section we will give an alternative description of these scalars, using the minimal polynomials $\Psi'_m$ of $2\cos(\pi/m)$ over $\ZZ$. 
First we will show that a pair of distinct roots of $\Psi'_m$ never sum to $0$.

\begin{lem} \label{lem:psiprimebezout}
Suppose $m>1$ is odd. 
We have
\begin{equation}
\left(\prod_{1 \neq d|m-2} \Psi'_d(-x)\right)\left(\prod_{1 \neq d|m}\Psi'_d(x)\right)+\left(\prod_{1 \neq d|m-2} \Psi'_d(x)\right)\left(\prod_{1 \neq d|m}\Psi'_d(-x)\right) = \pm 2 \text{.} \label{eq:psiprimebezout}
\end{equation}
\end{lem}

\begin{proof}
Let $C_n=2T_n(x/2)$ be a family of rescaled Chebyshev polynomials of the first kind. 
When $n=2k+1$ is odd we have
\begin{equation*}
\prod_{d|n} \Psi'_d = \frac{C_{2k+2}-C_{2k}}{C_{k+1}-C_k}
\end{equation*}
from the main result of \cite{wz-minpolcos}. 
Since $\Psi'_1=x+2$ we obtain
\begin{equation}
\prod_{1 \neq d|n} \Psi'_d = \frac{C_{2k+2}-C_{2k}}{(x+2)(C_{k+1}-C_k)} \text{.} \label{eq:psiprimeprod}
\end{equation}

We will need a few identities involving the rescaled Chebyshev polynomials. 
First, we note that $C_k$ has the same parity as $k$, i.e.~$C_k(-x)=(-1)^k C_k$ for any integer $k \geq 0$.
Second, we will need the product identity $C_j C_k = C_{j+k} + C_{|j-k|}$ for any integers $j,k \geq 0$. 
As a consequence, for any integer $k \geq 0$ we have
\begin{align}
(C_{k}-C_{k-1})(C_{k+1}+C_k) - (C_{k}+C_{k-1})(C_{k+1}-C_k) & =2C_k^2 - 2C_{k-1}C_{k+1} \nonumber \\
& =2C_{2k}+2C_0-2C_{2k}-2C_2  \label{eq:chebprod1} \\
& =2(4-x^2) \nonumber
\end{align}
and
\begin{align}
(C_k-C_{k-1})(C_k+C_{k-1}) & =C_k^2 - C_{k-1}^2 \nonumber \\
& =C_{2k}+C_0-(C_{2k-2}+C_0) \label{eq:chebprod2}\\
& =C_{2k}-C_{2k-2} \text{.} \nonumber
\end{align}

Now suppose $m=2n+1$. 
Using \eqref{eq:psiprimeprod} we can simplify the left-hand side of \eqref{eq:psiprimebezout}:
\begin{multline*}
\left(\prod_{1 \neq d|m-2} \Psi'_d(-x)\right)\left(\prod_{1 \neq d|m}\Psi'_d(x)\right)+\left(\prod_{1 \neq d|m-2} \Psi'_d(x)\right)\left(\prod_{1 \neq d|m}\Psi'_d(-x)\right) \\
\begin{aligned}
& = \frac{(C_{2n}(-x)-C_{2n-2}(-x))(C_{2n+2}-C_{2n})}{(-x+2)(C_{n}(-x)-C_{n-1}(-x))(x+2)(C_{n+1}-C_n)} \\
& \quad \quad {} + \frac{(C_{2n}-C_{2n-2})(C_{2n+2}(-x)-C_{2n}(-x))}{(x+2)(C_{n}-C_{n-1})(-x+2)(C_{n+1}(-x)-C_n(-x))} \\
%& =(-1)^n \frac{(C_{2n}-C_{2n-2})(C_{2n+2}-C_{2n})}{(4-x^2)(C_{n}+C_{n-1})(C_{n+1}-C_n)} \\
%& \quad \quad {} + (-1)^{n+1}\frac{(C_{2n}-C_{2n-2})(C_{2n+2}-C_{2n})}{(4-x^2)(C_{n}-C_{n-1})(C_{n+1}+C_n)} \\
& =(-1)^n \frac{(C_{2n}-C_{2n-2})(C_{2n+2}-C_{2n})}{(4-x^2)(C_{n}+C_{n-1})(C_{n+1}-C_n)} \\
& \quad \quad {} -(-1)^n\frac{(C_{2n}-C_{2n-2})(C_{2n+2}-C_{2n})}{(4-x^2)(C_{n}-C_{n-1})(C_{n+1}+C_n)} \text{.}% \\
\end{aligned}
\end{multline*}
Taking common denominators and applying \eqref{eq:chebprod1} and \eqref{eq:chebprod2} immediately gives the result:
\begin{multline*}
(-1)^n \frac{(C_{2n}-C_{2n-2})(C_{2n+2}-C_{2n})(C_{n}-C_{n-1})(C_{n+1}+C_n)}{(4-x^2)(C_{n}-C_{n-1})(C_n+C_{n-1})(C_{n+1}-C_n)(C_{n+1}+C_n)} \\
- (-1)^n\frac{(C_{2n}-C_{2n-2})(C_{2n+2}-C_{2n})(C_{n}+C_{n-1})(C_{n+1}-C_n)}{(4-x^2)(C_{n}-C_{n-1})(C_n+C_{n-1})(C_{n+1}-C_n)(C_{n+1}+C_n)} \\
\begin{aligned}
& =(-1)^n\left(\frac{2(C_{2n}-C_{2n-2})(C_{2n+2}-C_{2n})(4-x^2)}{(4-x^2)(C_{2n}-C_{2n-2})(C_{2n+2}-C_{2n})}\right) \\
& =2(-1)^n \qedhere
\end{aligned}
\end{multline*} 
\end{proof}

\begin{cor} \label{cor:psiprimerootsum}
Let $\Bbbk$ be a commutative integral domain, and suppose $m$ is a positive integer. 
If $r \in \Bbbk$ is a root of $\Psi'_m$, then $-r$ is a root as well if and only if $r=-r$.
\end{cor}

\begin{proof}
Lemma \ref{lem:psiprimebezout} gives a linear combination of $\Psi'_m(x)$ and $\Psi'_m(-x)$ which is equal to $\pm 2$. 
If both $r$ and $-r$ are roots of $\Psi'_m$, then they are roots of \eqref{eq:psiprimebezout} as well, so $\pm 2=0$ in $\Bbbk$, and therefore $r=-r$.
\end{proof}

\begin{thm} \label{thm:mminus1twocol}
Let $V$ be a realization of $(W,S)$ over $\Bbbk$.
Suppose $m=m_{st}>1$ is odd.
%Define two-colored quantum numbers $[2]_s=-\langle \alpha_s^\vee, \alpha_t\rangle$ and $[2]_t=-\langle \alpha_t^\vee, \alpha_s\rangle$. 
Then $\langle \alpha_s^\vee, \alpha_t \rangle \langle \alpha_t^\vee, \alpha_s\rangle$ is invertible and has a square root in $\Bbbk$. 
Moreover, if we fix the sign of this square root so that $\Psi'_m(\sqrt{\langle \alpha_s^\vee, \alpha_t\rangle \langle \alpha_t^\vee, \alpha_s\rangle})=0$ using Corollary \ref{cor:psiprimerootsum} then 
\begin{equation}
[m-1]_s = -\frac{\langle \alpha_s^\vee,\alpha_t\rangle}{\sqrt{\langle \alpha_s^\vee, \alpha_t\rangle \langle \alpha_t^\vee, \alpha_s\rangle}} \text{.} \label{eq:mminus1sqrt}
\end{equation}
%for some choice of square root.
%Moreover, the sign of the square root can be determined uniquely chosen so that $\Psi'_m(\sqrt{\langle \alpha_s^\vee, \alpha_t\rangle \langle \alpha_t^\vee, \alpha_s\rangle})=0$, where $\Psi'_m$ denotes the minimal polynomial of $2\cos(\pi/m)$ over $\ZZ$. 
\end{thm}

\begin{rem}
%As above let $V$ be a realization of $(W,S)$ over $\Bbbk$, and suppose $m=m_{st}>1$ is odd. 
Since $[m]_s=[m]_t=0$ from the definition of a realization, it is not hard to show that $[m-1]_s[m-1]_t=1$ (see e.g.~\cite[(6.11)]{ew-loccalc}).
Using this and \eqref{eq:2colqnumfrom1col} below it follows that
\begin{equation*}
[m-1]_s^2=\frac{\langle \alpha_s^\vee, \alpha_t\rangle}{\langle \alpha_t^\vee, \alpha_s \rangle} \text{.}
\end{equation*}
%using standard two-colored quantum identities. 
In other words, the real content of the theorem is in the determination of the \emph{sign} of $[m-1]_s$.
\end{rem}

\begin{proof}
%Consider the extension $A'=\ZZ[x,x_s,x_t] / (x^2-x_s x_t)$ of $A=\ZZ[x_s,x_t]$.
%In $A'$ we can define generic two-colored quantum numbers as in Definition \ref{defn:2colqnum}.
Recall the (generic) one-colored quantum numbers, defined as univariate polynomials in a variable $x$ by setting $[0]=0$, $[1]=1$, $[2]=x$, and inductively defining $[2][n]=[n+1]+[n-1]$.
These are essentially Chebyshev polynomials of the second kind (e.g.~\cite[\S 3]{2colJW}), i.e.
\begin{equation*}
[n](2\cos \theta)=U_{n-1}(\cos \theta)=\frac{\sin n\theta}{\sin \theta} \text{.} %\label{eq:1coltrig}
\end{equation*}
%Both \eqref{eq:2colqnumfrom1col} and \eqref{eq:1coltrig} apply 
The parity of $[n]$ (as a polynomial) is the opposite of the parity of $n$. 
In particular, when $n$ is even, $[2]$ divides $[n]$.

We also recall the \defnemph{cyclotomic parts}
\begin{equation*}
\Theta_n = \prod_{\substack{1 \leq k<n\\ (k,n)=1}} \left(x-2\cos \frac{k\pi}{n}\right)
\end{equation*}
from \cite[\S 3]{2colJW}, for which $\Theta_n(x)=\Psi_n(x^2)$ when $n>2$. 
For odd $n>1$ it is not hard to see that $\Psi'_n(x)\Psi'_n(-x)=\Theta_n$, as every root $2\cos(k\pi/n)$ of $\Theta_n$ is either a root of $\Psi'_n(x)$ or of $\Psi'_n(-x)$ depending on whether $k$ is odd or even.

The polynomial $[m-1]-1$ has $x=2\cos(\pi/m)$ as a root. 
Therefore $\Psi'_m$ divides $[m-1]-1$, and similarly $\Psi'_m(-x)$ divides $[m-1]+1$. 
This means that $\Psi'_m(x)\Psi'_m(-x)=\Psi_m(x^2)$ divides $[m-1]^2-1$. 
Moreover, as both polynomials are even, their quotient is even too.
%We can rephrase this entirely in the ring $\ZZ[x^2]/(\Psi_m(x^2))$.
Thus $u=\frac{[m-1]}{[2]} \in \ZZ[x^2]/(\Psi_m(x^2))$ is a root of the polynomial equation
\begin{equation}
u^2 x^2 - 1 \in (\ZZ[x^2]/(\Psi_m(x^2)))[u] \text{.} \label{eq:invsqroot}
\end{equation}
We can determine the sign of this root by observing that $[m-1][2]=[m-1]x$ is always a root of $\Psi'_m$ modulo $\Psi_m(x^2)$.

Now we use the relationship between one-colored and two-colored quantum numbers.
It can be shown (see e.g.~\cite[(6)]{2colJW}) that
\begin{equation}
\begin{aligned}
[n]_{s}& =[n](\sqrt{x_{s} x_{t}})=[n]_{t} & & \text{if $n$ is odd,} \\
\frac{[n]_{s}}{[2]_{s}}& =\left(\frac{[n]}{[2]}\right)(\sqrt{x_{s} x_{t}})=\frac{[n]_{t}}{[2]_{t}} & & \text{if $n$ is even.}
\end{aligned} \label{eq:2colqnumfrom1col}
\end{equation}
in $A=\ZZ[x_s,x_t]$.
Working in the quotient ring $B=\ZZ[x_s,x_t]/(\Psi_{m}(x_s x_t))$ it follows from \eqref{eq:invsqroot} that
\begin{equation*}
u=\frac{[m-1]_s}{[2]_s}=\frac{[m-1]}{[2]}(\sqrt{x_s x_t})
\end{equation*}
is a root of the equation $u^2 x_s x_t - 1 \in B[t]$.
In other words $x_s x_t$ has an invertible square root in $B$ which is a root of $\Psi'_m$; its reciprocal is equal to $\frac{[m-1]_s}{[2]_s}$.
We complete the proof by substituting $x_s=-\langle \alpha_s^\vee, \alpha_t\rangle$ and $x_t=-\langle \alpha_t^\vee, \alpha_s\rangle$. 
%
%In the quotient ring $\ZZ[x]/(\Psi'_m)$ we have
%\begin{equation*}
%[m-1] \equiv 1 \pmod{\Psi'_m} % \label{eq:mminusone}
%\end{equation*}
%since the difference between the left- and right-hand sides has $x=2\cos(\pi/m)$ as a root. 
%We can rewrite this as an equation only involving even polynomials. 
%Since $m$ is odd, we have $\Psi'_m(x)\Psi'_m(-x)=\Psi_m(x^2)$, because $\Psi'_m$ is irreducible and both $\Psi'_m$ and $\Psi_m$ are monic polynomials of degree $\varphi(m)/2$ with a common root $2\cos(\pi/m)$.
%Thus $\ZZ[x^2]/(\Psi'_m(x))=\ZZ[x^2]/(\Psi_m(x^2))$ makes sense as a quotient ring. 
%This means that $\frac{[m-1]}{[2]} \in \ZZ[x^2]/(\Psi_m(x^2))$ is a root of the equation
%\begin{equation*}
%t^2 x^2 - 1 \in (\ZZ[x^2]/(\Psi'_m))[t] \text{.}
%\end{equation*}
%
%Now we use the relationship between two-colored and one-colored quantum numbers. 
%It can be shown (see e.g.~\cite[(6)]{2colJW}) that
%\begin{equation*}
%\begin{aligned}
%[n]_{s}& =[n](\sqrt{x_{s} x_{t}})=[n]_{t} & & \text{if $n$ is odd,} \\
%\frac{[n]_{s}}{[2]_{s}}& =\left(\frac{[n]}{[2]}\right)(\sqrt{x_{s} x_{t}})=\frac{[n]_{t}}{[2]_{t}} & & \text{if $n$ is even.}
%\end{aligned} %\label{eq:2colqnumfrom1col}
%\end{equation*}
%in $A=\ZZ[x_s,x_t]$.
%Therefore, in the quotient $B=\ZZ[x_s,x_t]/(\Psi_{m}(x_s x_t))$ it follows that
%\begin{equation*}
%\frac{[m-1]_s}{[2]_s}=\frac{[m-1]}{[2]}(\sqrt{x_s x_t})
%\end{equation*}
%is a solution of the equation $t^2 x_s x_t - 1 \in B[t]$, i.e.~$x_s x_t$ has an invertible square root in $B$, which is a root of $\Psi'_m$. 
%We complete the proof by substituting $x_s=-\langle \alpha_s^\vee, \alpha_t\rangle$ and $x_t=-\langle \alpha_t^\vee, \alpha_s\rangle$. 
%
%If we interpret $x=2\cos(\pi/m)$, we observe that
%\begin{equation*}
%[m-1] = \frac{\sin \frac{(m-1)\pi}{m}}{\sin \frac{\pi}{m}} = 1
%\end{equation*}
%so 
%\begin{equation*}
%[m-1]_s = \frac{[m-1][2]_s}{[2]} = \frac{x_s}{x} \in \left\{\pm \frac{x_s}{\sqrt{x_s x_t}}\right\} \text{.}
%\end{equation*}
%We complete the proof by substituting $x_s=-\langle \alpha_s^\vee, \alpha_t\rangle$ and $x_t=-\langle \alpha_t^\vee, \alpha_s\rangle$. shows that \eqref{eq:mminus1sqrt} holds.
%
%We can determine the sign of the square root if only one of
%\begin{equation*}
%\pm\frac{-\langle x_s^\vee, x_t\rangle}{\sqrt{\frac{\langle x_s^\vee, x_t\rangle}{\langle x_t^\vee, x_s\rangle}}} = \pm \sqrt{\langle \alpha_s^\vee, \alpha_t\rangle \langle \alpha_t^\vee, \alpha_s \rangle}
%\end{equation*}
%satisfies $\Psi_m'$. 
%In other words, if we can show that $\Psi'_m$ \emph{never} has a pair of roots which sum to $0$ (except when $2=0$ in $\Bbbk$) then the result follows.
\end{proof}

\begin{cor} \label{cor:balancedcriterion}
Let $V$ be a realization of $(W,S)$ over $\Bbbk$.
Then $V$ is balanced if and only if $\Psi'_{m_{st}}(-\langle \alpha_s^\vee, \alpha_t \rangle)=0$ for all distinct $s,t \in S$ such that $m_{st}$ is odd. 
\end{cor}



\section{Demazure surjectivity and cinched Hecke categories}

%Here we apply the results of the previous sections to the construction of diagrammatic Hecke categories. 

%\subsection{Demazure realizations}

%Let $V$ be a $\Bbbk$-realization of $(W,S)$, and let $R=\Sym(V)$ be the symmetric algebra of $V$ with $\deg V=2$. 
%For any $s \in S$ write $R^s$ for the subring of invariants with respect to $s$. 
%The following maps, called \defnemph{Demazure operators}, play a vital role in the theory of Soergel bimodules.
%
%\begin{prop}[{\cite[Exercise 5.46]{emtw}}] \label{prop:demazureopdefn}
%There is a unique homogeneous $R^s$-bimodule homomorphism $\partial_s: R \rightarrow R^s(-2)$ (i.e.~a graded homomorphism in degree $-2$) which satisfies the following properties.
%\begin{enumerate}
%\item For $v \in V$, $\partial_s(v)=\langle \alpha_s, v\rangle$.
%
%\item The twisted Leibniz rule holds: for $f,g \in R$, we have $\partial_s(fg)=\partial_s(f)g+s(f)\partial_s(g)$. 
%\end{enumerate}
%Moreover, the following identities hold: 
%\begin{align*}
%\ident-s& =\alpha_s\partial_s \text{,} & \partial_s^2& =0 \text{,} \\
%s \circ \partial_s& =\partial_s\text{,} & \partial_s \circ s& =-\partial_s \text{.}
%\end{align*}
%\end{prop}
%
%\begin{proof}
%We will show that the linear map $\partial_s : V \rightarrow \Bbbk$ defined by $\partial_s=\alpha_s^\vee$ extends to a homogeneous linear map $\partial_s : R \rightarrow R(-2)$ via the twisted Leibniz rule. 
%To show that it is well-defined, we need to check that the twisted Leibniz rule respects associativity and commutativity. 
%More precisely, we must show that for any $f,g,h \in R$, applying the twisted Leibniz rule to $\partial_s(f(gh))$ and $\partial_s((fg)h)$ gives the same result, and similarly that applying the twisted Leibniz rule to $\partial_s(fg)$ and $\partial_s(gf)$ gives the same result. 
%
%Respecting associativity is completely straightforward:
%\begin{gather*}
%\partial_s(f(gh))=\partial_s(f)gh+s(f)\partial_s(gh)=\partial_s(f)gh+s(f)\partial_s(g)h+s(f)s(g)\partial_s(h) \\
%\partial_s((fg)h)=\partial_s(fg)h+s(fg)\partial_s(h)=\partial_s(f)gh+s(f)\partial_s(g)h+s(f)s(g)\partial_s(h) \text{.}
%\end{gather*}
%
%For commutativity, we will first prove that $f-s(f)=\partial_s(f)\alpha_s$ by induction on degree.
%By definition this holds in degrees $0$ and $2$. 
%So suppose by induction it holds for all polynomials of degree less than $k$ for some $k>2$. 
%Any monomial of degree $k$ is a product $fg$ of two monomials of degree less than $k$.
%Using the induction hypothesis we have
%\begin{align*}
%fg-s(fg)& =fg-s(f)g+s(f)g-s(f)s(g) \\
%& =\partial_s(f)\alpha_s g+s(f)\partial_s(g)\alpha_s \\
%& =\partial_s(fg)\alpha_s
%\end{align*}
%so the result holds for all monomials of degree $k$, and by linearity, for all polynomials of degree $k$, and hence for all polynomials by induction.
%
%It follows that $\partial_s$ respects commutativity, since
%\begin{align*}
%\partial_s(fg)& =\partial_s(f)g+s(f)\partial_s(g) \\
%& =\partial_s(f)g-\partial_s(f)s(g)+\partial_s(f)s(g)+s(f)\partial_s(g) \\
%& =\partial_s(f)\partial_s(g)\alpha_s+\partial_s(f)s(g)+s(f)\partial_s(g) \\
%& =f\partial_s(g)+\partial_s(f)s(g) \\
%& =\partial_s(gf) \text{.}
%\end{align*}
%The remaining identities follow by induction on degree in a similar manner.
%\end{proof}
%
%\begin{rem}
%In the literature the Demazure operators are usually defined by the formula
%\begin{equation*}
%\partial_s(f)=\frac{f-s(f)}{\alpha_s}
%\end{equation*}
%for $f \in R$. 
%It is not immediately clear from this formula that $\partial_s$ is well defined. 
%If $\alpha_s \neq 0$, then one can use induction on degree to show that $f-s(f)$ is always a (well-defined) multiple of $\alpha_s$. 
%If in addition $\alpha_s^\vee:V \rightarrow \Bbbk$ is surjective, then it can be shown that $R \cong R^s \oplus R^s(-2)$ as $R^s$-bimodules, and that $\partial_s$ is the projection onto the second direct summand \cite[\S 3.3]{EW-SoergelCalc}
%%A standard argument to show that $\partial_s$ is well defined proceeds by assuming that $\alpha_s^\vee: V \rightarrow \Bbbk$ is surjective, obtaining a decomposition $R \cong R^s \oplus R^s(-2)$ and then proving that $\partial_s$ is the projection onto the second direct summand \cite[\S 3.3]{EW-SoergelCalc}. 
%%Alternatively, one can assume that $\alpha_s \neq 0$ and then induct on degree to show that $f-s(f)$ is always a (well-defined) multiple of $\alpha_s$. 
%Proposition \ref{prop:demazureopdefn} above shows that neither of these assumptions is really necessary. 
%\end{rem}


An extremely common assumption on realizations in the literature is Demazure surjectivity, which we recall (and slightly generalize) below.

\begin{defn} \label{defn:demazuresurj}
Let $V$ be a realization of $(W,S)$ over $\Bbbk$.
\begin{itemize}
\item We say $V$ is \defnemph{coroot surjective} if $\alpha_s^\vee : V \rightarrow \Bbbk$ is surjective for all $s \in S$.

\item We say $V$ is \defnemph{root surjective} if $\alpha_s : V^\ast \rightarrow \Bbbk$ is surjective for all $s \in S$.

\item We say $V$ is \defnemph{Demazure surjective} if it is both coroot surjective and root surjective.
\end{itemize}
\end{defn}


Demazure surjectivity is often conflated with coroot surjectivity in the literature. 
For example, it is sometimes claimed that Demazure surjectivity is \emph{necessary} for the Demazure operators $\{\partial_s\}$ to be surjective, or that every realization embeds inside a Demazure surjective realization (e.g.~\cite[\S 2.3]{gjw-calcpcanbasis}). 
These are both incorrect as written, but the corresponding statements for coroot surjective realizations are true.

%\begin{lem}
%The Demazure operators are surjective if and only if the underlying realization is coroot surjective.
%Moreover, every realization embeds inside a coroot surjective realization. 
%\end{lem}
%
%\begin{proof}
%If $\partial_s: R \rightarrow R^s(-2)$ is surjective, then it is surjective in degree $2$, i.e.~$(\partial_s)_2=\alpha_s^\vee:V \rightarrow \Bbbk$ is surjective. 
%Conversely, if $\delta \in V$ such that $\partial_s(\delta)=1$ then for any $f \in R^s$ we have $\partial_s(\delta f)=f$ by the twisted Leibniz rule.
%
%For the second statement, define a new realization $V'=V \oplus \bigoplus_{s \in S} \Bbbk f_s$ by setting $\langle \alpha_s^\vee, f_t \rangle = \delta_{s,t}$. 
%By construction $V'$ is coroot surjective, and it is immediate that the natural map $V \hookrightarrow V'$ is a morphism of realizations.
%\end{proof}

\begin{exam}
Let $\Bbbk=\mathbb{F}_2$ and $(W,S)$ be the $A_1$ Coxeter group. 
We consider two different realization structures on the one-dimensional vector space $V=\Bbbk \varepsilon$. 
Write $\varepsilon^\ast \in V^\ast$ for the dual basis vector, i.e.~$\varepsilon^\ast(\varepsilon)=1$.
\begin{enumerate}
\item Let $\alpha_s=0 \in V$ and $\alpha_s^\vee=\varepsilon^\ast \in V^\ast$. 
Clearly $V$ is coroot surjective but not root surjective, so $V$ is not Demazure surjective.
We also observe that $\partial_s(\varepsilon)=\alpha_s^\vee(\varepsilon)=1$, so the Demazure operator $\partial_s$ is surjective.

\item Let $\alpha_s=0 \in V$ and $\alpha_s^\vee=0 \in V^\ast$. 
Now $V$ is neither coroot surjective nor root surjective. 
If $V'$ is another realization and $\phi:V \rightarrow V'$ is a morphism of realizations, then $\alpha_s'=\phi(\alpha_s)=\phi(0)=0$. 
In particular $V'$ cannot be Demazure surjective. 
\end{enumerate}
\end{exam}


%\begin{rem}
%
%In fact, $\partial_s$ can be defined in the absence of \emph{any} assumptions (including that $\alpha_s \neq 0$!). 
%First we define $\partial_s$ on $V$ by setting $\partial_s(v)= \langle \alpha_s^\vee, v \rangle$ for any $v \in V$. 
%Then we extend to a map $\partial_s: \Sym(V) \rightarrow \Sym(V)$ using the twisted Leibniz rule $\partial_s(fg)=\partial_s(f)g+s(f)\partial_s(g)$. 
%To show that this is well-defined, we need to check associativity and commutativity. 
%More precisely, we must show that applying the twisted Leibniz rule to $\partial_s(f(gh))$ and $\partial_s((fg)h)$ gives the same result, and similarly that applying the twisted Leibniz rule to $\partial_s(fg)$ and $\partial_s(gf)$ gives the same result. 
%Associativity is completely straightforward. 
%For commutativity, we can proceed by induction using the auxiliary claim that $f-s(f)=\partial_s(f)\alpha_s$.
%
%To check well-definedness here, we need to check associativity and commutativity, i.e.~that evaluating $\partial_s(fgh)$ in two different ways gives
%(At this stage it is necessary to check associativity, i.e.~that evaluating $\partial_s(fgh)$ in two different ways gives the same result.)
%Then we show that $\partial_s$ descends to a map $\partial_s : \Sym(V) \rightarrow \Sym(V)$. 
%One way to do this is by induction on degree, showing that 
%
% $\partial_s: \Sym(V)_i \rightarrow \Sym(V)_{i-1}$
%It is easy to show (by induction on degree) that $f-s(f)=\partial_s(f)\alpha_s$, and thus 
%\end{rem}

%Of the conditions above, coroot surjectivity is arguably the most natural, as it ensures that the Demazure operators
%\footnote{In fact, many treatments use coroot surjectivity to \emph{define} the Demazure operators. This is not necessary --- the formula $\partial_s(f)=(f-s(f))/\alpha_s$ is always well defined!}
%$\{\partial_s\}_{s \in S}$ are surjective.

Let $V$ be a $\Bbbk$-realization of $(W,S)$, and let $R=\Sym(V)$ be the symmetric algebra of $V$ with $\deg V=2$. 
For $s \in S$ write $R^s$ for the subring of invariants with respect to $s$. 
We will say a little more about the relationship between the various surjectivity conditions and different incarnations of the Hecke category. 

In the classical (i.e.~non-diagrammatic) theory of Soergel bimodules, surjectivity of the Demazure operator $\partial_s$ gives $R^s \subset R$ the structure of a Frobenius extension, which can be used to show that the Bott--Samelson bimodule $B_s=R \otimes_{R^s} R(1)$ is a Frobenius algebra object (see e.g.~\cite[Lemma~8.27]{emtw}).
In the diagrammatic approach the corresponding object $B_s$ is automatically given a Frobenius algebra object structure by diagrammatic relations. 
Here coroot surjectivity is only necessary to show that the tensor square $B_s^{\otimes 2}$ splits as a direct sum $B_s(1) \oplus B_s(-1)$ (see e.g.~\cite[(5.14)]{EW-SoergelCalc}), i.e.~the simplest case of Soergel's categorification theorem \cite[Theorem~11.1]{emtw}.

By contrast, root surjectivity on its own is rather strong. 
A simpler assumption (namely, that the roots are non-zero) is enough to ensure that the localization functor is well defined \cite{ew-loccalc}.
This combined with coroot surjectivity implies that Soergel's categorification theorem holds.
%For example, to define the localization functor on the Hecke category it is enough to assume that the roots are merely non-zero. 
Demazure surjectivity is a stronger assumption, but it is convenient because it ensures that Soergel's categorification theorem also holds for the \emph{dual} realization.

\begin{rem}
One notable situation where the roots being surjective (as opposed to just non-zero) matters arises from the left monodromy action on the relevant \defnemph{mixed derived categories} \cite[\S 2.3]{amrw-koszulduality}. 
It can be shown that the monodromy maps for the generating tilting objects are surjective if and only if the realization is root surjective \cite[Example~4.7.4]{amrw-freemonodromic}. 
As this action is closely connected to Koszul duality for Hecke categories (which exchanges the underlying realization with its dual), it is perhaps unsurprising that root surjectivity is meaningful here. 
\end{rem}

%For these reasons, when working with the Elias--Williamson diagrammatic Hecke category is it usually enough to assume coroot surjectivity and that all the roots are non-zero.

%Demazure surjectivity is often considered essential in the theory of Hecke categories, as it ensures that the diagrammatic Hecke category actually categorifies the Hecke algebra. 

In fact \emph{none} of the surjectivity conditions in Definition \ref{defn:demazuresurj} are necessary to construct a diagrammatic categorification of the Hecke algebra.

Let $\mathcal{D}_{\rm BS}$ denote the Elias--Williamson diagrammatic category for an arbitrary $\Bbbk$-realization $V$ of a Coxeter system $(W,S)$ (introduced in \cite{EW-SoergelCalc}). 
It turns out that a slightly modified version of $\mathcal{D}_{\rm BS}$ \emph{always} categorifies the Hecke algebra, and thus deserves the moniker ``Hecke category''. 
(This approach was first outlined without proof in \cite[Remark~1.15]{bdhn-hsp}.) 
%In what follows we will show that a slightly modified version of $\mathcal{D}_{\rm BS}$ is always well behaved, even in the absence 
In what follows we will \emph{not} call $\mathcal{D}_{\rm BS}$ ``the diagrammatic Bott--Samelson category'' or its Karoubi envelope $\mathcal{D}$ ``the diagrammatic Hecke category'', to avoid potential confusion.
%if one is willing to add an additional relation to the definition of the diagrammatic Hecke category. 
%Below we propose a modified definition of the (diagrammatic) Hecke category which is slightly more well behaved.
The proposed replacement for $\mathcal{D}_{\rm BS}$ involves adding a single extra diagrammatic relation.

\begin{defn}
The \defnemph{(cinched) diagrammatic Bott--Samelson category} $\mathcal{H}_{\rm BS}$ is the quotient of $\mathcal{D}_{\rm BS}$ with respect to the additional diagrammatic (i.e.~monoidal) ``cinching relation'' 
\begin{equation}
\begin{minipage}{1.5cm}\begin{tikzpicture}[scale=1.000]
\draw[densely dotted, rounded corners] (-0.5,0) rectangle (1,1.5) ;
\draw[red,line width=0.08cm](0.5,0)--++(90:1.5);
\draw[red,line width=0.08cm](0,0)--++(90:1.5);   \end{tikzpicture}
  \end{minipage}
  \;=\;
   \begin{minipage}{1.5cm}\begin{tikzpicture}[scale=1.000]
\draw[densely dotted, rounded corners] (-0.5,0) rectangle (1,1.5) ;
%\draw[red,line width=0.08cm](0.25,0.4)--++(90:1.5-0.8);
%%\draw[red,line width=0.08cm](0.5,0)--++(90:1.5);
%\draw[red,line width=0.08cm](0,0) to [out=90,in=180]
%(0.25,0.4)  to [out=0,in=90]
%(0.5,0);
%
\draw[red,line width=0.08cm](0,1.5) to [out=-90,in=180]
(0.5,1.5-0.45)   ;
 
 
  \draw[red,line width=0.08cm](0.5,0)--++(90:1.5);

\draw[red,line width=0.08cm] (0 ,0)--++(90:0.45) coordinate(hi);
\fill[red] (hi) circle (3.5pt);
  \end{tikzpicture}
  \end{minipage}
  %%
\;  + \;  \begin{minipage}{1.5cm}\begin{tikzpicture}[scale=1.000,yscale=-1.000]
\draw[densely dotted, rounded corners] (-0.5,0) rectangle (1,1.5) ;
%\draw[red,line width=0.08cm](0.25,0.4)--++(90:1.5-0.8);
%%\draw[red,line width=0.08cm](0.5,0)--++(90:1.5);
%\draw[red,line width=0.08cm](0,0) to [out=90,in=180]
%(0.25,0.4)  to [out=0,in=90]
%(0.5,0);
%
\draw[red,line width=0.08cm](0,1.5) to [out=-90,in=180]
(0.5,1.5-0.45)   ;
 
 
  \draw[red,line width=0.08cm](0.5,0)--++(90:1.5);

\draw[red,line width=0.08cm] (0 ,0)--++(90:0.45) coordinate(hi);
\fill[red] (hi) circle (3.5pt);
  \end{tikzpicture}
  \end{minipage}
  %%
\; - 
 \;
 \begin{minipage}{1.5cm}\begin{tikzpicture}[scale=1.000]
\draw[densely dotted, rounded corners] (-0.5,0) rectangle (1,1.5) ;
%\draw[red,line width=0.08cm](0.25,0.4)--++(90:1.5-0.8);
%%\draw[red,line width=0.08cm](0.5,0)--++(90:1.5);
%\draw[red,line width=0.08cm](0,0) to [out=90,in=180]
%(0.25,0.4)  to [out=0,in=90]
%(0.5,0);
%
\draw[red,line width=0.08cm](0,1.5) to [out=-90,in=180]
(0.5,1.5-0.45)   ;
\draw[red,line width=0.08cm](0,0) to [out=90,in=180]
(0.5,0.45)   ;


\draw[red,line width=0.08cm](0.5,0)--++(90:1.5);
\draw[red,line width=0.08cm] (0-0.2 ,0.75)--++(90:0.3) coordinate(hi);
\fill[red] (hi) circle (3.5pt);
\draw[red,line width=0.08cm] (0-0.2 ,0.75)--++(-90:0.3) coordinate(hi);
\fill[red] (hi) circle (3.5pt);
  \end{tikzpicture}
  \end{minipage} \label{eq:cinching}
\end{equation}
for each color.
We call the Karoubi envelope of the additive graded closure of $\mathcal{H}_{\rm BS}$ the \defnemph{(cinched) diagrammatic Hecke category} $\mathcal{H}$.
\end{defn}


The cinching relation \eqref{eq:cinching} first appeared as an independent relation in \cite[(R4)]{bdhn-hsp} along with informal versions of many of the results below. 
Adding a dot to the right strand gives an equivalent form of the cinching relation:
%\missingfigure{dottedcinch}
\begin{equation}
\begin{minipage}{1.5cm}
\begin{tikzpicture}[scale=1.000]
\draw[densely dotted, rounded corners] (-0.5,0) rectangle (1,1.5) ;
\draw[red,line width=0.08cm](0.5,0)--++(90:0.45) coordinate (dot);
\fill[red] (dot) circle (3.5pt);
\draw[red,line width=0.08cm](0,0)--++(90:1.5);   \end{tikzpicture}
  \end{minipage}
  \;=\;
   \begin{minipage}{1.5cm}\begin{tikzpicture}[scale=1.000]
\draw[densely dotted, rounded corners] (-0.5,0) rectangle (1,1.5) ;
%\draw[red,line width=0.08cm](0.25,0.4)--++(90:1.5-0.8);
%%\draw[red,line width=0.08cm](0.5,0)--++(90:1.5);
%\draw[red,line width=0.08cm](0,0) to [out=90,in=180]
%(0.25,0.4)  to [out=0,in=90]
%(0.5,0);
%
%\draw[red,line width=0.08cm](0,1.5) to [out=-90,in=180]
%(0.5,1.5-0.45)   ;
 
 
  \draw[red,line width=0.08cm](0.5,0)--++(90:1.5);

\draw[red,line width=0.08cm] (0 ,0)--++(90:0.45) coordinate(hi);
\fill[red] (hi) circle (3.5pt);
  \end{tikzpicture}
  \end{minipage}
  %%
\;  + \;  \begin{minipage}{1.5cm}\begin{tikzpicture}[scale=1.000,yscale=-1.000]
\draw[densely dotted, rounded corners] (-0.5,0) rectangle (1,1.5) ;
%\draw[red,line width=0.08cm](0.25,0.4)--++(90:1.5-0.8);
%%\draw[red,line width=0.08cm](0.5,0)--++(90:1.5);
%\draw[red,line width=0.08cm](0,0) to [out=90,in=180]
%(0.25,0.4)  to [out=0,in=90]
%(0.5,0);
%
\draw[red,line width=0.08cm](0,1.5) arc (180:360:0.25);
%(0.5,1.5)   ;
 
 
%  \draw[red,line width=0.08cm](0.5,0)--++(90:1.5);

\draw[red,line width=0.08cm] (0.25,0)--++(90:0.45) coordinate(hi);
\fill[red] (hi) circle (3.5pt);
  \end{tikzpicture}
  \end{minipage}
  %%
\; - 
 \;
 \begin{minipage}{1.5cm}\begin{tikzpicture}[scale=1.000]
\draw[densely dotted, rounded corners] (-0.5,0) rectangle (1,1.5) ;
%\draw[red,line width=0.08cm](0.25,0.4)--++(90:1.5-0.8);
%%\draw[red,line width=0.08cm](0.5,0)--++(90:1.5);
%\draw[red,line width=0.08cm](0,0) to [out=90,in=180]
%(0.25,0.4)  to [out=0,in=90]
%(0.5,0);
%
%\draw[red,line width=0.08cm](0,1.5) to [out=-90,in=180]
%(0.5,1.5-0.45)   ;
\draw[red,line width=0.08cm](0,0) to [out=90,in=180]
(0.5,0.45)   ;


\draw[red,line width=0.08cm](0.5,0)--++(90:1.5);
\draw[red,line width=0.08cm] (0-0.2 ,0.75)--++(90:0.3) coordinate(hi);
\fill[red] (hi) circle (3.5pt);
\draw[red,line width=0.08cm] (0-0.2 ,0.75)--++(-90:0.3) coordinate(hi);
\fill[red] (hi) circle (3.5pt);
  \end{tikzpicture}
  \end{minipage} \label{eq:dottedcinch}
\end{equation}

\begin{rem} \hfill
\begin{enumerate}
\item We have included the term ``cinched'' here only to avoid potential confusion with the original Elias--Williamson construction. 
Once we have shown that $\mathcal{H}_{\rm BS}$ is better behaved, and in particular \emph{always} categorifies the Hecke algebra, it is reasonable to redefine the terms ``diagrammatic Bott--Samelson category'' and ``diagrammatic Hecke category'' so that they are always ``cinched'' (i.e.~always include the cinching relation) (cf.~\cite{bowman-book}).

\item Many results for $\mathcal{H}_{\rm BS}$ carry over without change. 
For example, it is easy to see that the construction of the cinched Bott--Samelson category is functorial in the realization, and that localization is still well defined whenever the roots of the underlying realization are all non-zero. 
%The primary application of the latter is to prove the linear independence of the light leaves maps. 
%In fact this holds regardless of the realization. 
\end{enumerate}
\end{rem}

Recall the \defnemph{light leaves} and \defnemph{double leaves} constructions for $\mathcal{D}_{\rm BS}$ \cite[\S 6.1, \S 6.3]{EW-SoergelCalc}. 
These diagrammatic constructions carry over identically to the cinched Hecke category. %$\mathcal{H_{\rm BS}$, and similarly yields a basis.
The main result in this section is the following. %that the resulting maps similarly yield bases for various $\Hom$-spaces in $\mathcal{H}_{\rm BS}$. 
%As promised, the main result for in this section is that Soergel's categorification theorem always holds for cinched Hecke categories.


\begin{thm}[{cf.~\cite[Theorem 6.11, Proposition 7.6]{EW-SoergelCalc}}] \label{thm:LLcinched} \hfill
%Let $\underline{x}, \underline{y}$ be expressions, and let $\underline{w}$ be a reduced expression for $w \in W$.
%\begin{enumerate}
%\item The light leaves maps $\{\mathrm{LL}_{\underline{x},\mathbf{e}} : \mathbf{e} \text{ is a subsequence of $\underline{x}$ expressing } w\}$ form a basis for $\Hom_{\mathcal{H}_{\rm BS}}(B_{\underline{x}},B_{\underline{w}})/I_{\underline{w}}$, where $I_{\underline{w}}$ denotes the ideal of lower terms (defined in \cite[\S 7.2]{EW-SoergelCalc}).
%
%\item The double leaves maps $\mathbb{LL}_{\underline{x},\underline{y}}$ form a cellular basis for $\Hom_{\mathcal{H}_{\rm BS}}(B_{\underline{x}},B_{\underline{y}})$.
%\end{enumerate}

The light leaves and double leaves constructions yield bases for the $\Hom$-spaces of the cinched Bott--Samelson category $\mathcal{H}_{\rm BS}$.
\end{thm}

As in $\mathcal{D}_{\rm BS}$ this immediately implies the following.

\begin{cor}[{cf.~\cite[Theorem 11.1, Assumption 11.25]{emtw}}] \label{cor:cinchedSoergelcatthm} \hfill

Suppose $\Bbbk$ is a henselian\footnote{E.g.~a field, or a complete local ring.} local ring. 
Then Soergel's categorification theorem holds for the cinched Hecke category $\mathcal{H}$.
\end{cor}



%
%\begin{lem}
%Let $\phi: V \rightarrow V'$ be a $\Bbbk$-morphism of realizations, and write $R$ and $R'$ for the corresponding symmetric algebras. 
%The map $\phi$ extends to a $\Bbbk$-linear monoidal functor $\phi:\mathcal{H}_{\rm BS} \rightarrow \mathcal{H}'_{\rm BS}$ between the respective Bott--Samelson categories, mapping generators to generators.
%This functor descends to a functor $\phi:\mathcal{H}^c_{BS} \rightarrow (\mathcal{H}'_{BS})^c$ on the cinched Bott--Samelson categories. 
%Finally, the image of each $\Hom$-space is naturally a $\phi(R)$-bimodule, and 
%\begin{align*}
%\Hom_{\mathcal{H}'_{\rm BS}}(B_{\underline{x}},B_{\underline{y}}) & = R' \otimes_{\phi(R)} \phi(\Hom_{\mathcal{H}_{\rm BS}}(B_{\underline{x}},B_{\underline{y}}))\text{,} \\
%\Hom_{(\mathcal{H}'_{\rm BS})^c}(B_{\underline{x}},B_{\underline{y}}) & = R' \otimes_{\phi(R)} \phi(\Hom_{\mathcal{H}^c_{\rm BS}}(B_{\underline{x}},B_{\underline{y}}))
%\end{align*}
%for any expressions $\underline{x}$ and $\underline{y}$. 
%\end{lem}
%%More precisely,  any morphism of $\Bbbk$-realizations $V \rightarrow V'$ extends to a monoidal functor $\mathcal{H}_{\rm BS} \rightarrow \mathcal{H}'_{\rm BS}$ (resp.~$\mathcal{H}_{\rm BS}^c \rightarrow (\mathcal{H}'_{\rm BS})^c$) between the two (cinched) Bott--Samelson categories mapping generators to generators in the obvious way. 
%%This functor is $\Bbbk$-linear
%
%\begin{proof}
%It is straightforward to check that the functor $\phi$ is well defined. 
%For the final statement on $\Hom$-spaces, note that the polynomial forcing relation \cite[(5.2)]{EW-SoergelCalc} means that each $\Hom$-space in $\mathcal{H}'_{\rm BS}$ is spanned (as a left $R'$-module) by the Soergel diagrams without polynomials. 
%As such diagrams are in the image of $\phi$ by construction, this proves the result.
%\end{proof}
%
%In general 
%
%\begin{exam}
%Let $\Bbbk=\mathbb{F}_2$ and $(W,S)$ be the $A_1$ Coxeter group.
%It is easy to check that $V=\Bbbk$ with $\alpha_s=1$ and $\alpha_s^\vee=0$ defines a $\Bbbk$-realization of $(W,S)$. 
%The realization $V$ is not coroot surjective; it embeds into the coroot surjective realization $V'=\Bbbk^2$ defined by $\alpha_s=(1,0)$, $\alpha_s^\vee=(0,1)$.
%Thus
%\begin{equation*}
%\Hom_{\mathcal{H}'_{\rm BS}}(B_{\underline{ss}}, B_{\underline{s}}) = R' \otimes_{\phi(R)} \phi(\Hom_{\mathcal{H}_{\rm BS}}(B_{\underline{ss},B_{\underline{s}}})) \text{,}
%\end{equation*}
%where $\phi: V\rightarrow V'$ denotes this natural embedding.
%In particular the presence of $\phi$ on the right-hand side is essential, since \eqref{eq:dottedcinch} does \emph{not} hold in $\mathcal{H}_{\rm BS}$ but does hold in $\mathcal{H}'_{\rm BS}$.
%\end{exam}

\begin{proof}[Proof of Theorem \ref{thm:LLcinched}]
We follow the same strategy as in the proof of \cite[Proposition~7.6]{EW-SoergelCalc}, by first showing that the light leaves maps are linearly independent, and then showing that they span the relevant $\Hom$-space modulo lower terms. 

If we assume that the roots in $V$ are all non-zero then the usual localization argument carries over to $\mathcal{H}_{\rm BS}$ to prove linear independence. 
For completeness we will show that this assumption is unnecessary.

\begin{clm} \label{clm:LLlinindep}
The light leaves maps in $\mathcal{H}_{\rm BS}$ are linearly independent.
\end{clm}

%The usual argument to prove this uses localization, which works just as well in the cinched Bott--Samelson category, as long as all the roots in $V$ are non-zero.
%which relies on the (often unremarked) assumption that the all the roots are nonzero. 

\begin{proof}
%We will prove this for $\mathcal{D}_{\rm BS}$; the argument for $\mathcal{H}_{\rm BS}$ is identical.
Let $I=\{t \in S : \alpha_t=0\}$. 
In $V$ we observe that $\langle \alpha^\vee_s,\alpha_t \rangle=0$ whenever $t \in I$, so in particular $w(\alpha_t)=\alpha_t$ for any $t \in I$.
Define a new realization
\begin{equation*}
V'=V \oplus \bigoplus_{s \in I} \Bbbk \varepsilon_s
\end{equation*}
by setting
\begin{align*}
\alpha'_s& =\alpha_s & &\text{if $s \notin I$,} \\
\alpha'_s& =\varepsilon_s & &\text{if $s \in I$,} \\
\langle (\alpha'_s)^\vee, v \rangle& =\langle \alpha_s^\vee, v\rangle & &  \text{if  $s \in S$ and $v \in V$,} \\
\langle (\alpha'_s)^\vee, \varepsilon_t\rangle & =0 & &  \text{if $s \in S$ and $t \in I$.}
\end{align*}
%Let $A$ be the Cartan matrix of the underlying realization $V$, and let $V_{A,\Bbbk}$ be the universal realization with respect to this Cartan matrix. 
%Let $U$ be the span of the roots in $V$, and let $f_1,\dotsc,f_r$ be a basis for the complement of $U$ inside $V$. 
%We define a new realization
%\begin{equation*}
%V' = V_{A,\Bbbk} \oplus \left(\bigoplus_{i=1}^r \Bbbk' f'_i\right)
%\end{equation*}
%by setting 
%\begin{equation*}
%\langle \alpha_s^\vee, f'_i \rangle = \langle \alpha_s^\vee, f_i\rangle \qquad \qquad \qquad \text{ for all $s \in S$ and $1 \leq i \leq r$.}
%\end{equation*}
Let $\mathcal{H}'_{\rm BS}$ denote the cinched Bott--Samelson category constructed from $V'$, and let $K=\bigoplus_{s \in I} \Bbbk \alpha'_s \leq V'$.
Due to the form of the diagrammatic relations, it is evident that $\mathcal{H}_{\rm BS}$ is the quotient of $\mathcal{H}'_{\rm BS}$ modulo any Soergel diagram with a term in $K$ in some region. 
Any such term in $K$ is fixed by $W$, so we can push this term all the way to the left using the polynomial forcing relation \cite[(5.2)]{EW-SoergelCalc} without introducing any extra error terms.
This shows that $\mathcal{H}_{\rm BS} = R \otimes_{R'} \mathcal{H}_{\rm BS}$, where $R$ and $R'$ denote the symmetric algebras of $V$ and $V'$ respectively.  
%We will show that $\mathcal{H}_{\rm BS}=R \otimes_{R'} \mathcal{H}'_{\rm BS}$, 
%We first observe that $\mathcal{H}_{\rm BS}$ is equivalent to the quotient of $\mathcal{H}'_{\rm BS}$
Now observe that the light leaves maps in $\mathcal{H}'_{\rm BS}$ are linearly independent by the usual localization argument, as the roots of $V'$ are all non-zero.
In each relevant $\Hom$-space, the light leaves maps for $\mathcal{H}'_{\rm BS}$ span a free left $R'$-submodule. 
Changing scalars to $R$ preserves freeness, so the corresponding light leaves maps are linearly independent in $\mathcal{H}_{\rm BS}$ too.
\end{proof}

To show that these maps also span the relevant $\Hom$-space modulo lower terms, the proof is identical to the usual argument in \cite[\S 7]{EW-SoergelCalc}, noting that the only use of coroot surjectivity there is in the proof of \cite[Claim~7.10]{EW-SoergelCalc}. 
In $\mathcal{H}_{\rm BS}$ it is enough to prove a slightly weaker version of this statement. 
See \cite[\S\S 6--7]{EW-SoergelCalc} for a description of the notation (and the definition of ``lower terms'').

\begin{clm}
Let $\underline{x}$ be an expression in $S$, and let $\mathbf{e}$ be a subsequence of $\underline{x}$ which expresses some $w \in W$.
Suppose $s \in S$ such that $ws<w$. 
Then %$\mathrm{LL}_{\underline{x},\mathbf{e}} \otimes {\rm dot}_s$ 
%\missingfigure{LLdot}
\begin{equation*}
\mathrm{LL}_{\underline{x},\mathbf{e}} \otimes {\rm dot}_s = 
\begin{tikzpicture}[scale=0.4,baseline=(std)] %,every node/.append style={transform shape, font={\Huge}}]
\def\strings{6}
\pgfmathsetmacro{\stringgap}{1.0/\strings}
%
\node (std) at (0,0) [draw, rectangle, minimum width=2.4cm, minimum height=0.6cm] {$\mathrm{LL}_{\underline{x},\mathbf{e}}$};
%
\foreach \i in {1,...,\strings}
{
    \pgfmathsetmacro{\xoffset}{(\i-0.5)*\stringgap}
    \path ($(std.north west)!\xoffset!(std.north east)$) node (topb\i) {};
    \path (topb\i) ++(0,1.5cm) node (topa\i) {};
    \path ($(std.south west)!\xoffset!(std.south east)$) node (bota\i) {};
    \path (bota\i) %++(0,-1cm) node (botb\i) {}
                   ++(0,-1.5cm) node (botc\i) {};
}
\pgfmathsetmacro{\dotoffset}{(\strings+0.75)*\stringgap}
\path ($(std.west)!\dotoffset!(std.east)$) node[dot,fill=red] (dot) {};
\path (dot.center |- botc6.center) node (botdot) {};
% bot edges
\foreach \i in {1,...,6}
{
    \pgfmathsetmacro{\xoffset}{(\i-0.5)*\stringgap}
    \draw (bota\i.center) edge[string] (botc\i.center);
}
\draw (bota3.center) edge[string,red] (botc3.center);
\draw (dot.center) edge[string,red] (botdot.center);
% top edges
\foreach \i in {2,...,5}
{
    \pgfmathsetmacro{\xoffset}{(\i-0.5)*\stringgap}
    \draw (topa\i.center) edge[string] (topb\i.center);
}
\draw (topa5.center) edge[string,red] (topb5.center);
\end{tikzpicture}
\end{equation*}
can be written as a linear combination of light leaves maps with polynomials in any region modulo lower terms.
\end{clm}

\begin{proof}
Apply the modified form of the cinching relation \eqref{eq:dottedcinch}:
%\missingfigure{LLcinched}
\begin{equation*}
\begin{tikzpicture}[scale=0.4,baseline=(std)] %,every node/.append style={transform shape, font={\Huge}}]
\def\strings{6}
\pgfmathsetmacro{\stringgap}{1.0/\strings}
%
\node (std) at (0,0) [draw, rectangle, minimum width=2.4cm, minimum height=0.6cm] {$\mathrm{LL}_{\underline{x},\mathbf{e}}$};
%
\foreach \i in {1,...,\strings}
{
    \pgfmathsetmacro{\xoffset}{(\i-0.5)*\stringgap}
    \path ($(std.north west)!\xoffset!(std.north east)$) node (topb\i) {};
    \path (topb\i) ++(0,3.5cm) node (topa\i) {};
    \path ($(std.south west)!\xoffset!(std.south east)$) node (bota\i) {};
    \path (bota\i) %++(0,-1cm) node (botb\i) {}
                   ++(0,-1.5cm) node (botc\i) {};
}
\pgfmathsetmacro{\dotoffset}{(\strings+0.75)*\stringgap}
\path ($(std.west)!\dotoffset!(std.east)$) node[dot,fill=red] (dot) {};
\path (dot.center |- botc6.center) node (botdot) {};
% bot edges
\foreach \i in {1,...,6}
{
    \pgfmathsetmacro{\xoffset}{(\i-0.5)*\stringgap}
    \draw (bota\i.center) edge[string] (botc\i.center);
}
\draw (bota3.center) edge[string,red] (botc3.center);
\draw (dot.center) edge[string,red] (botdot.center);
% top edges
\foreach \i in {2,...,5}
{
    \pgfmathsetmacro{\xoffset}{(\i-0.5)*\stringgap}
    \draw (topa\i.center) edge[string] (topb\i.center);
}
\draw (topa5.center) edge[string,red] (topb5.center);
\end{tikzpicture}
=
\begin{tikzpicture}[scale=0.4,baseline=(std)] %,every node/.append style={transform shape, font={\Huge}}]
\def\strings{6}
\pgfmathsetmacro{\stringgap}{1.0/\strings}
%
\node (std) at (0,0) [draw, rectangle, minimum width=2.4cm, minimum height=0.6cm] {$\mathrm{LL}_{\underline{x},\mathbf{e}}$};
%
\foreach \i in {1,...,\strings}
{
    \pgfmathsetmacro{\xoffset}{(\i-0.5)*\stringgap}
    \path ($(std.north west)!\xoffset!(std.north east)$) node (topb\i) {};
    \path (topb\i) ++(0,3.5cm) node (topa\i) {};
    \path ($(std.south west)!\xoffset!(std.south east)$) node (bota\i) {};
    \path (bota\i) %++(0,-1cm) node (botb\i) {}
                   ++(0,-1.5cm) node (botc\i) {};
}
\pgfmathsetmacro{\dotoffset}{(\strings+0.75)*\stringgap}
\path ($(std.west)!\dotoffset!(std.east)$) node (dot) {};
\path (dot.center |- botc6.center) node (botdot) {};
\path ($(topa5.center)!0.25!(topb5.center)$) node (turn1) {};
\path (turn1.center -| botdot.center) node (turn2) {};
\path ($(turn1.center)!0.5!(turn2.center)$) node (turn) {};
\path ($(topa5.center)!0.75!(topb5.center)$) node[dot,fill=red] (dotnew) {};
% bot edges
\foreach \i in {1,...,6}
{
    \pgfmathsetmacro{\xoffset}{(\i-0.5)*\stringgap}
    \draw (bota\i.center) edge[string] (botc\i.center);
}
\draw (bota3.center) edge[string,red] (botc3.center);
\draw (dot.center) edge[string,red] (botdot.center);
% top edges
\foreach \i in {2,...,4}
{
    \pgfmathsetmacro{\xoffset}{(\i-0.5)*\stringgap}
    \draw (topa\i.center) edge[string] (topb\i.center);
}
\draw[corner, string, red] (topa5.center) |- (turn.center) -| (botdot.center);
\draw[string,red] (dotnew.center) edge (topb5.center);
\end{tikzpicture}
+
\begin{tikzpicture}[scale=0.4,baseline=(std)] %,every node/.append style={transform shape, font={\Huge}}]
\def\strings{6}
\pgfmathsetmacro{\stringgap}{1.0/\strings}
%
\node (std) at (0,0) [draw, rectangle, minimum width=2.4cm, minimum height=0.6cm] {$\mathrm{LL}_{\underline{x},\mathbf{e}}$};
%
\foreach \i in {1,...,\strings}
{
    \pgfmathsetmacro{\xoffset}{(\i-0.5)*\stringgap}
    \path ($(std.north west)!\xoffset!(std.north east)$) node (topb\i) {};
    \path (topb\i) ++(0,3.5cm) node (topa\i) {};
    \path ($(std.south west)!\xoffset!(std.south east)$) node (bota\i) {};
    \path (bota\i) %++(0,-1cm) node (botb\i) {}
                   ++(0,-1.5cm) node (botc\i) {};
}
\pgfmathsetmacro{\dotoffset}{(\strings+0.75)*\stringgap}
\path ($(std.west)!\dotoffset!(std.east)$) node (dot) {};
\path (dot.center |- botc6.center) node (botdot) {};
\path ($(topa5.center)!0.25!(topb5.center)$) node[dot,fill=red] (dotnew) {};
\path ($(topa5.center)!0.75!(topb5.center)$) node (turn1) {};
\path (turn1.center -| botdot.center) node (turn2) {};
\path ($(turn1.center)!0.5!(turn2.center)$) node (turn) {};
% bot edges
\foreach \i in {1,...,6}
{
    \pgfmathsetmacro{\xoffset}{(\i-0.5)*\stringgap}
    \draw (bota\i.center) edge[string] (botc\i.center);
}
\draw (bota3.center) edge[string,red] (botc3.center);
%\draw (dot.center) edge[string,red] (botdot.center);
% top edges
\foreach \i in {2,...,4}
{
    \pgfmathsetmacro{\xoffset}{(\i-0.5)*\stringgap}
    \draw (topa\i.center) edge[string] (topb\i.center);
}
\draw[corner, string, red] (topb5.center) |- (turn.center) -| (botdot.center);
\draw[string,red] (dotnew.center) edge (topa5.center);
\end{tikzpicture}
-
\begin{tikzpicture}[scale=0.4,baseline=(std)] %,every node/.append style={transform shape, font={\Huge}}]
\def\strings{6}
\pgfmathsetmacro{\stringgap}{1.0/\strings}
%
\node (std) at (0,0) [draw, rectangle, minimum width=2.4cm, minimum height=0.6cm] {$\mathrm{LL}_{\underline{x},\mathbf{e}}$};
%
\foreach \i in {1,...,\strings}
{
    \pgfmathsetmacro{\xoffset}{(\i-0.5)*\stringgap}
    \path ($(std.north west)!\xoffset!(std.north east)$) node (topb\i) {};
    \path (topb\i) ++(0,3.5cm) node (topa\i) {};
    \path ($(std.south west)!\xoffset!(std.south east)$) node (bota\i) {};
    \path (bota\i) %++(0,-1cm) node (botb\i) {}
                   ++(0,-1.5cm) node (botc\i) {};
}
\pgfmathsetmacro{\dotoffset}{(\strings+0.75)*\stringgap}
\path ($(std.west)!\dotoffset!(std.east)$) node (dot) {};
\path (dot.center |- botc6.center) node (botdot) {};
\path ($(topa5.center)!0.25!(topb5.center)$) node (turn1) {};
\path (turn1.center -| botdot.center) node (turn2) {};
\path ($(turn1.center)!0.5!(turn2.center)$) node (turn) {};
\path ($(topa5.center)!0.75!(topb5.center)$) node (dotnew) {};
\path ($(topa5.center)!0.375!(topb5.center)$) node[dot,fill=red] (barbella) {};
\path ($(topa5.center)!0.625!(topb5.center)$) node[dot,fill=red] (barbellb) {};
% bot edges
\foreach \i in {1,...,6}
{
    \pgfmathsetmacro{\xoffset}{(\i-0.5)*\stringgap}
    \draw (bota\i.center) edge[string] (botc\i.center);
}
\draw (bota3.center) edge[string,red] (botc3.center);
\draw (dot.center) edge[string,red] (botdot.center);
% top edges
\foreach \i in {2,...,4}
{
    \pgfmathsetmacro{\xoffset}{(\i-0.5)*\stringgap}
    \draw (topa\i.center) edge[string] (topb\i.center);
}
\draw[corner, string, red] (topa5.center) |- (turn.center) -| (botdot.center);
\draw[corner,string,red] (topb5.center) |- (dotnew.center -| botdot.center);
\draw[string,red] (barbella.center) edge (barbellb.center);
\end{tikzpicture}
\end{equation*}
The second term on the right-hand side vanishes modulo lower terms as it factors through $B_{\underline{ws}}$.
The last term is already a light leaves map with an extra polynomial $\alpha_s$. 
So it is enough to reduce the first term. 
By the Jones--Wenzl relation \cite[(5.7)]{EW-SoergelCalc} dots ``propagate'' down through braids (see also \cite[(4.4)]{matrixrecursion}). 
Thus we can push the dot through any rex move until it either reaches the bottom of the diagram or it reaches a trivalent vertex.
In the first case we obtain a new light leaves map with a $\mathrm{U}0$-strand in place of a $\mathrm{U}1$-strand.
In the second case, we obtain a ``birdcage''
%\missingfigure{birdcage}
\begin{equation*}
\begin{tikzpicture}[scale=0.4,baseline=(origin.center)]
\coordinate (origin) at (0,0) {};
\path (origin) ++(0,-1.5cm) coordinate (bot1) {}
               ++(1cm,0) coordinate (bot2) {}
               ++(1cm,0) coordinate (bot3) {}
               ++(1.5cm,0) node (ellipses) {$\dotso$}
               ++(1.5cm,0) coordinate (bot4) {};
\path (origin) ++(0,3cm) coordinate (top1) {};
\path ($(bot1.center)!0.33!(top1.center)$) coordinate (along2);
\path ($(bot1.center)!0.5!(top1.center)$) coordinate (along3);
\path ($(bot1.center)!0.66!(top1.center)$) coordinate (along4);
\path ($(bot1.center)!0.85!(top1.center)$) node[dot,fill=red] (dot) {};
\draw[string,red] (bot1.center) edge (dot.center);
\draw[corner,string,red] (bot2.center) |- (along2.center)
                  (bot3.center) |- (along3.center)
                  (bot4.center) |- (along4.center);
\end{tikzpicture}
=
\begin{tikzpicture}[scale=0.4,baseline=(origin.center)]
\coordinate (origin) at (0,0) {};
\path (origin) ++(0,-1.5cm) coordinate (bot1) {}
               ++(1cm,0) coordinate (bot2) {}
               ++(1cm,0) coordinate (bot3) {}
               ++(1.5cm,0) node (ellipses) {$\dotso$}
               ++(1.5cm,0) coordinate (bot4) {};
\path (origin) ++(0,3cm) coordinate (top1) {};
\path ($(bot1.center)!0.33!(top1.center)$) coordinate (along2);
\path ($(bot1.center)!0.5!(top1.center)$) coordinate (along3);
\path ($(bot1.center)!0.66!(top1.center)$) ++(1cm,0) coordinate (turn);
%\path ($(bot1.center)!0.85!(top1.center)$) node[dot,fill=red] (dot) {};
%\draw[string,red] (bot1.center) edge (dot.center);
\draw[corner,string,red] (bot2.center) |- (along2.center)
                  (bot3.center) |- (along3.center)
                  (bot4.center) |- (turn.center) -| (bot1.center);
\end{tikzpicture}
\end{equation*}
which also forms part of a light leaves map: the first strand of the birdcage is $\mathrm{U}1$, the last strand of the birdcage is $\mathrm{D}1$, and any intermediate strands are $\mathrm{D}0$.
\end{proof}

This shows that the light leaves maps for a basis for certain $\Hom$-spaces modulo lower terms.
Applying the same argument as \cite[\S 7.3]{EW-SoergelCalc} shows that the double leaves form bases for the entire $\Hom$-spaces of $\mathcal{H}_{\rm BS}$.
\end{proof}

\begin{rem}
The situation for diagrammatic \emph{singular} Hecke categories is completely different. 
In this setting the invariant rings $R^I$ for all \defnemph{finitary} $I \subseteq S$ (i.e.~for which $W_I$ is finite) play a key role. 
For all finitary $I \subset J$, the diagrammatic relations themselves posit the existence of a Demazure operator $\partial_J^I : R^I \rightarrow R^J$ which gives $R^J \subset R^I$ a Frobenius extension structure (see e.g.~\cite[\S 24.2.1]{emtw}. %by the diagrammatic relations alone (see e.g.~\cite[(24.14)]{emtw}). 
In particular, these Demazure operators must all be surjective for the diagrammatic singular Hecke category to be well defined. 

Unlike in the regular setting this condition naturally places restrictions on the base ring $\Bbbk$ and the Cartan matrix $A$. 
To see this, suppose $V$ is a realization for which the Demazure operators $\{\partial_J^I : I \subset J \text{ finitary}\}$ are all surjective. 
There is a natural morphism $\phi:V \rightarrow V_{A,{\rm sc}}$ by Lemma \ref{lem:adscuniversality}, which extends to a morphism $\phi:R \rightarrow R_{A,{\rm sc}}$ between the corresponding symmetric algebras.
One can check that $\partial_J^I = (\partial_J^I)_{A,{\rm sc}} \circ \phi$, where $(\partial_J^I)_{A,{\rm sc}}$ denotes the corresponding Demazure operator for $V_{A,{\rm sc}}$. 
This means that $V_{A,{\rm sc}}$ must have surjective Demazure operators too. 
The problem of determining precisely when this can happen dates back to Demazure's original paper introducing Demazure operators \cite{demazure}!
For the standard Cartan matrices of Weyl groups, he carefully showed that $2$ must be invertible outside of types $A$ and $C$, $3$ must be invertible in types $E$ and $F$, and $5$ must be invertible in type $E_8$.
\end{rem}

\section{Applications}





We conclude this paper with some applications of cinched Bott--Samelson and cinched Hecke categories. 
Many of these results have known analogues for $\mathcal{D}_{\rm BS}$ and $\mathcal{D}$ (with the added assumption of coroot surjectivity); the primary utility of the cinching construction is to simplify the proofs. 
As above $V$ is a $\Bbbk$-realization from which the categories $\mathcal{H}_{\rm BS}$, $\mathcal{H}$, $\mathcal{D}_{\rm BS}$, and $\mathcal{D}$ have been constructed.

The following is a version of \cite[Remark~3.19]{EW-SoergelCalc} for cinched Bott--Samelson categories.

\begin{lem} \label{lem:changerlz}
Let $V \rightarrow V'$ be a morphism of realizations, with symmetric algebras $R$ and $R'$ respectively.
The induced monoidal functor $\mathcal{H}_{\rm BS} \rightarrow \mathcal{H}'_{\rm BS}$ gives rise to an isomorphism $\mathcal{H}_{\rm BS} \cong R' \otimes_{R} \mathcal{H}_{\rm BS}$ of $\Bbbk$-linear categories, i.e.
\begin{equation*}
\Hom_{\mathcal{H}'_{\rm BS}}(B_{\underline{x}},B_{\underline{y}}) \cong R' \otimes_R \Hom_{\mathcal{H}_{\rm BS}}(B_{\underline{x}},B_{\underline{y}})
\end{equation*}
for any expressions $\underline{x},\underline{y}$ in $S$.
%The functor decategorifies to an isomorphism $[\mathcal{H}] \cong [\mathcal{H}']$ of split Grothendieck rings.
If $\Bbbk$ is a henselian local ring, then the Karoubi envelope of this functor maps indecomposable objects to indecomposable objects. %, and the corresponding canonical bases of the Hecke algebra coincide.
\end{lem}

\begin{proof}
The induced monoidal functor maps light leaves to light leaves. 
Since the double leaves form bases for the $\Hom$-spaces this proves the first part.
It also shows that the local intersection forms (see e.g.~\cite[Definition~11.74]{emtw}) for the two categories coincide. 
The second part follows because local intersection forms determine the decompositions of the Bott--Samelson objects into direct sums of indecomposable objects \cite[Corollary~11.75]{emtw}
\end{proof}

We can use this to say a little more about the relationship between the cinched Bott--Samelson category $\mathcal{H}_{\rm BS}$ and the Elias--Williamson diagrammatic category $\mathcal{D}_{\rm BS}$.

\begin{prop}
%Let $V$ be a realization of a Coxeter system $(W,S)$, and let $\mathcal{H}_{\rm BS}$ be the corresponding cinched diagrammatic Bott--Samelson category. 
%Then $\mathcal{H}^c$ always categorifies the Hecke algebra of $(W,S)$. 
%Moreover, if $V$ is Demazure surjective, then $\mathcal{H}^c$ is equivalent to the usual diagrammatic Hecke category $\mathcal{H}$ constructed from $V$.
%\begin{enumerate}
If $V$ is coroot surjective, then $\mathcal{H}_{\rm BS}$ is equivalent to the Elias--Williamson diagrammatic category $\mathcal{D}_{\rm BS}$.
More generally, let $V \hookrightarrow V'$ be an embedding of $V$ into a coroot surjective realization $V'$, which extends to an embedding $R \hookrightarrow R'$ of their symmetric algebras. 
Let $\mathcal{D}'_{\rm BS}$ denote the Elias--Williamson diagrammatic category constructed from $V'$. 
Then there is an embedding $\mathcal{H}_{\rm BS} \rightarrow \mathcal{D}'_{\rm BS}$ which induces an isomorphism $\mathcal{D}'_{\rm BS} \cong R' \otimes_R \mathcal{H}_{\rm BS}$. 
%(where $R=\Sym(V)$ and $R'=\Sym(V')$). 
%i.e.
%\begin{equation*}
%\Hom_{\mathcal{D}'_{\rm BS}}(B_{\underline{x}},B_{\underline{y}}) \cong R' \otimes_R \Hom_{\mathcal{H}_{\rm BS}}(B_{\underline{x}},B_{\underline{y}}) \text{,}
%\end{equation*}
%for any expressions $\underline{x}$ and $\underline{y}$ in $S$.
\end{prop}

In other words, the cinched Bott--Samelson category $\mathcal{H}_{\rm BS}$ is equivalent to the usual Elias--Williamson diagrammatic category $\mathcal{D}_{\rm BS}$ whenever the latter is well behaved.
Moreover $\mathcal{H}_{\rm BS}$ is always a monoidal subcategory of a well-behaved Elias--Williamson diagrammatic category $\mathcal{D}'_{\rm BS}$. 

\begin{proof}
If $V$ is coroot surjective, then the cinching relation already holds in $\mathcal{D}_{\rm BS}$ \cite[(5.14)--(5.15)]{EW-SoergelCalc}. 

For general $V$, let $V'$ be a coroot surjective realization with an embedding $V \hookrightarrow V'$, and let $\mathcal{H}'_{\rm BS}$ denote the corresponding cinched Bott--Samelson category.
The composition of the induced functor $\mathcal{H}_{\rm BS} \rightarrow \mathcal{H}'_{\rm BS}$ with the equivalence $\mathcal{H}'_{\rm BS} \xrightarrow{\sim} \mathcal{D}'_{\rm BS}$ gives the $\Hom$-formula above.
\end{proof}

%\begin{rem}
%Koszul duality for the diagrammatic Elias--Williamson category $\mathcal{D}$ has been generalized from the original setting of affine Coxeter groups to arbitrary Coxeter groups \cite{rv-koszulduality}. 
%In all cases the proofs rely on Demazure surjectivity in a crucial way.
%It is an open question whether Koszul duality holds for cinched Hecke categories in complete generality.
%\end{rem}

%The following is probably not new (or even dependent on the cinched construction) but is so useful that it deserves stating in complete generality.
Next, we recall that a consequence of Soergel's categorification theorem is that the indecomposable objects give rise to a \defnemph{canonical basis} of the Hecke algebra. 
For example, the \defnemph{$p$-canonical} or \defnemph{$p$-Kazhdan--Lusztig basis} arises from the Hecke category for a Cartan realization of a finite or affine Weyl group over a field $\Bbbk$ of characteristic $p$. 
In principle such canonical bases depend on the realization, and indeed there are numerous examples which showcase dependency on the characteristic $p$ or the Cartan matrix $A$. 
In fact, these essentially turn out to be the only things the canonical basis can depend on. 
This has been observed before for the $p$-canonical basis \cite[Remark~10.17(2)]{rw-tilting}. 
The only potential difficulty in formulating a more general result is describing precisely what it means for two realizations over different base rings to have ``the same'' Cartan matrix.

\begin{prop}
Suppose $\Bbbk$ is a henselian local ring with residue field $\mathbb{F}$. 
Let $A$ denote the Cartan matrix of the realization $V$, and let $\mathbb{F}'$ be the subfield of $\mathbb{F}$ generated by its prime subfield and the images of the entries $a_{st}$ of $A$ in $\mathbb{F}$. 
The canonical basis of the Hecke algebra arising from $\mathcal{H}$ depends only on $\mathbb{F}'$ and image of $A$ in $\mathbb{F}'$. 
\end{prop}

\begin{proof}
%Let $A$ be the Cartan matrix for $V$, and let $V' \leq V$ denote the root subrealization of $V$ (i.e.~the image of the unique morphism from the universal realization $V_{A,\Bbbk}$ to $V$). 
If $B$ is an indecomposable object in $\mathcal{H}$, then its image in $\mathbb{F} \otimes_{\Bbbk} \mathcal{H}$ is also indecomposable by the usual idempotent lifting argument.
%Recall the following fundamental property of henselian local rings: if $B$ is a $\Bbbk$-algebra of finite $\Bbbk$-rank, then the idempotents of $B$ can be lifted from those of $\mathbb{F} \otimes_{\Bbbk} B$. 
%Applying this to $B=\End_{\mathcal{H}_{\rm BS}}^0(B_{\underline{x}})$ for some reduced expression $\underline{x}$ shows that the Bott--Samelson object $B_{\underline{x}}$ has identical decompositions in $\mathcal{H}$ and $\mathbb{F} \otimes_{\Bbbk} \mathcal{H}$. 
The latter is the cinched Hecke category for the $\mathbb{F}$-realization $\mathbb{F} \otimes_{\Bbbk} V$, so without loss of generality we may assume that $\Bbbk$ is a field. 

Let $\Bbbk'$ be the subfield of $\Bbbk$ generated by the prime subfield and the images of the entries of the Cartan matrix, and let $V'$ denote the restriction of $V$ to a $\Bbbk'$-module. 
Clearly $V'$ is a $\Bbbk'$-realization, and for the corresponding cinched Bott--Samelson category $\mathcal{H}_{\rm BS}'$ we have $\mathcal{H}_{\rm BS}=\Bbbk \otimes_{\Bbbk'} \mathcal{H}'_{\rm BS}$, and similarly for the local intersection forms of the two categories. 
Recall that the decomposition of each Bott--Samelson object into a direct sum of indecomposable objects is determined by the ranks of these forms. 
As the rank of a matrix does not change after field extension, this shows that indecomposable objects in $\mathcal{H}'$ remain indecomposable in $\mathcal{H}$.
So without loss of generality we may further assume that $\Bbbk=\Bbbk'$.

To complete the proof we consider the morphism $V_{A,{\rm ad}} \rightarrow V$, which gives rise to a monoidal functor $\mathcal{H}_{A,{\rm ad}} \rightarrow \mathcal{H}$ by Lemma \ref{lem:changerlz}.
This functor maps indecomposable objects to indecomposable objects, so the canonical bases arising from both cinched Hecke categories must coincide.
% where $\mathcal{H}_{A,\Bbbk}$ be the cinched Hecke category constructed from the universal realization $V_{A,\Bbbk}$.
\end{proof}

%As an immediate corollary is folklore, as evidenced by the terminology ``$p$-canonical basis''; see \cite[Remark~10.17(2)]{rw-tilting} for a proof in the Kac--Moody case.
%
%\begin{cor}
%Suppose $\Bbbk$ is a field of characteristic $p \geq 0$, and $A$ is a Cartan matrix defined over $\ZZ$. 
%Let $\mathcal{H}$ be a cinched Hecke algebra for a $\Bbbk$-realization $V$ whose
%The Hecke algebra canonical basis arising from a cinched Hecke category over $\Bbbk$ depends only on its Cartan matrix and the characteristic of $\Bbbk$. 
%\end{cor}



Finally we apply our cinched constructions to the parabolic anti-spherical setting, to show that any realization can be used to categorify the anti-spherical module of the Hecke algebra. 
%Let us assume in what follows that $\Bbbk$ is a henselian local ring. 

Fix some subset $I \subseteq S$. 
We write $W_I$ to denote the subgroup generated by $I$, and $\prescript{I}{}{W}$ for the set of minimal length representatives for the right cosets $W_I \backslash W$. 
An expression $\underline{x}$ in $S$ is called an \defnemph{$I$-sequence} if it starts with some element $s \in I$. 
We define the \defnemph{(cinched) anti-spherical Bott--Samelson category $\mathcal{N}_{\rm BS}$} to be the quotient of $\mathcal{H}_{\rm BS}$ by the ideal generated by $\{B_{\underline{x}} : \underline{x} \text{ is an $I$-sequence}\}$.
%Now let $V$ be a realization of $(W,S)$, and let $\mathcal{H}$ be the associated Hecke category. 
%Recall that the indecomposable objects $\{B_x\}$ of $\mathcal{H}$ are indexed by elements of $W$. 
%The \defnemph{anti-spherical category $\mathcal{N}$} is defined to be the quotient of $\mathcal{H}$ by the ideal generated by $\{B_x : x \notin \prescript{I}{}{W}\}$. 
It is a right module category for $\mathcal{H}_{\rm BS}$, i.e.~there is a right monoidal action of $\mathcal{H}_{\rm BS}$ on $\mathcal{N}_{\rm BS}$. 
We call the Karoubi envelope of $\mathcal{N}_{\rm BS}$ the \defnemph{(cinched) anti-spherical category $\mathcal{N}$}. 

Recall that the light leaves construction for $\mathcal{H}_{\rm BS}$ generalizes to $\mathcal{N}_{\rm BS}$, with anti-spherical light leaves and double leaves \cite[Definition~5.1]{LW-antispher}.

%In fact, this is true regardless of the realization, as we prove below.

\begin{thm}[{cf.~\cite[Theorem~5.3]{LW-antispher}}] \label{thm:antispherLLbasis}
The anti-spherical light leaves and double leaves constructions yield bases for the $\Hom$-spaces of the cinched anti-spherical Bott--Samelson category $\mathcal{N}_{\rm BS}$.
\end{thm}

%\begin{thm}
%The category $\mathcal{N}$ categorifies the anti-spherical module of the Hecke algebra. 
%\end{thm}

Libedinsky--Williamson proved this when the underlying realization $V$ is Demazure surjective and satisfies the ``parabolic property'' \cite[\S 2.3]{LW-antispher}.

\begin{proof}
%It is enough to show that the anti-spherical light leaves defined in \cite[\S 5]{LW-antispher} form bases (over an appropriate polynomial ring) for the spaces
%\begin{equation*}
%\Hom_{\mathcal{N}^{\geq x}}^{\bullet}(N_{\underline{y}}, D_x)
%\end{equation*}
%for $x \in \prescript{I}{}{W}$ and $\underline{y}$ an expression. 
%Here $N_{\underline{y}}$ denotes the image in $\mathcal{N}$ of the Bott--Samelson bimodule $B_{\underline{y}}$, $D_x$ denotes the image of $B_x$, and 
%$\mathcal{N}^{\geq x}$
%%\begin{equation*}
%%\mathcal{N}^{\geq x}= \mathcal{N} / \langle D_w(m) : m \in \ZZ \text{,}\ w \ngeq x \rangle_{\oplus}
%%\end{equation*}
%is the quotient of $\mathcal{N}$ with respect to the full additive subcategory generated by (grade shifts of) the indecomposable objects $D_w$ for $w \ngeq x$. 
%
Let $U$ denote the root subrealization of $V$, and let $K$ be a complement for $U$ inside $V$. 
Choose a basis $f_1,\dotsc,f_r$ of $K$. % such that 
%\begin{equation*}
%\{\alpha_s : s \in S\} \cup \{f_i : 1 \leq i \leq r\}
%\end{equation*}
%span all of $V$. 
Let $\mathbb{O}$ be the subring of $\overline{\QQ}$ generated by the algebraic integers $4\cos^2(\pi/m_{st})$ for all $s,t \in S$ such that $m_{st}<\infty$.
By definition, for every $s,t \in S$ the product $\langle \alpha_s^\vee,\alpha_t\rangle \langle \alpha_t^\vee, \alpha_s \rangle$ in $\Bbbk$ can be lifted to $\mathbb{O}$.
Next we use $\mathbb{O}$ to define a ring
\begin{equation*}
\Bbbk'=\frac{\mathbb{O}[x_{s,t},y_{s,i}: s,t \in S,\ 1 \leq i \leq r]}{(x_{s,t}x_{t,s} - 4\cos^2(\pi/m_{st}) \text{ for all $s,t \in S$ with $m_{st}<\infty$})} \text{.}
\end{equation*}
Now we define a new $\Bbbk'$-realization
\begin{equation*}
V' = \left(\bigoplus_{s \in S} \Bbbk'\alpha'_s \right) \oplus \left(\bigoplus_{i=1}^r \Bbbk' f'_i\right)
\end{equation*}
with linearly independent roots $\{\alpha'_s : s \in S\}$ by setting 
\begin{align*}
\langle (\alpha'_s)^\vee, f'_i \rangle & = y_{s,i} \text{ for all $s \in S$ and $1 \leq i \leq r$,} \\
\langle (\alpha'_s)^\vee, \alpha'_t \rangle & =x_{s,t} \text{ for all $s,t \in S$.}
\end{align*}

Consider the specialization of $V'$ where $x_{s,t}=-2\cos(\pi/m_{st})$ for all $s,t \in S$ (suitably interpreting $\cos(\pi/\infty)=\cos 0=1$). 
The root subrealization of this specialization is equivalent to the geometric representation, as defined in \cite[\S 5.3]{humphreys}. 
As this representation is faithful, the original realization $V'$ must also be faithful, and therefore satisfies the parabolic property. 
Thus by \cite[Theorem~5.3]{LW-antispher} the anti-spherical light leaves form bases for the relevant $\Hom$-spaces in $\mathcal{N}'_{\rm BS}$, where $\mathcal{N}'_{\rm BS}$ is the anti-spherical Bott--Samelson category defined by $V'$. 
But $V'$ also specializes to $V$ in the obvious way, by setting $x_{s,t}=\langle \alpha_s^\vee,\alpha_t\rangle$ and $y_{s,i}=\langle \alpha_s^\vee,f_i\rangle$ for all $s,t \in S$ and $1 \leq i \leq r$. 
Since the bases for the $\Hom$-spaces remain bases after specialization, this proves the result.
\end{proof}

\begin{cor}[{cf.~\cite[Theorem~6.2]{LW-antispher}}] \label{cor:antispherSoergelcatthm}
Let $\Bbbk$ be a henselian local ring. 
The category $\mathcal{N}$ categorifies the anti-spherical module of the Hecke algebra. 
\end{cor}

\begin{rem}
A simpler version of Theorem \ref{thm:antispherLLbasis} first appeared in \cite[Theorem~1.10]{BHN-modularWeylKac}. 
It was used to show that the anti-spherical category for Cartan realizations and geometric realizations is well behaved in positive characteristic \cite[Example~1.11]{BHN-modularWeylKac}. 
A related idea to justify defining the $p$-canonical basis via $p$-modular systems was described in \cite[\S 3.3]{gjw-calcpcanbasis}.
\end{rem}

\begin{rem} \label{rem:locpres}
The linear independence argument in Claim \ref{clm:LLlinindep} is completely new, and applies to the Elias--Williamson diagrammatic category $\mathcal{D}_{\rm BS}$ with no change.
As written it is only relevant when $2$ vanishes in $\Bbbk$, but the same type of argument works in other settings where roots vanish, such as anti-spherical Hecke categories (as in Theorem \ref{thm:antispherLLbasis}) and cyclotomic Hecke categories (analogous to categories of Soergel modules). 

Indeed, it is worth understanding how to reinterpret localization in such settings.
Carrying over the notation from Claim \ref{clm:LLlinindep} let $\Lambda' : \mathcal{H}'_{\rm BS} \rightarrow (\Omega_{Q'} W)_{\oplus}$ be the localization functor for $\mathcal{H}'_{\rm BS}$, as defined in \cite{ew-loccalc}. 
Let $\mathcal{M}'$ denote the image of $\Lambda'$ inside $(\Omega_{Q'} W)_{\oplus}$, which is isomorphic to $\mathcal{H}'_{\rm BS}$ because localization is faithful.
It has a concrete presentation in terms of matrices with entries in the fraction field $Q'$ of $R'$. 
Now let 
\begin{equation}
\Lambda = R \otimes_{R'} \Lambda' : R \otimes_{R'} \mathcal{H}'_{\rm BS} \xrightarrow{\sim} R \otimes_{R'} \mathcal{M}'
\end{equation}
This gives a faithful presentation of $\mathcal{H}_{\rm BS}$ in terms of a matrix algebra in which the scalars have been changed. 
Note in particular that the target of $\Lambda$ does \emph{not} lie in $Q \otimes_{Q'}(\Omega_{Q'} W)_{\oplus} = (\Omega_{Q} W)_{\oplus}$, where $Q$ denotes the fraction field of $R$. 
In other words, we cannot interpret the image of $\Lambda$ as a matrix algebra with entries in $Q$. 
For example, if $\alpha_s=0$ in $R$ then
\begin{equation*}
1 \otimes 
\begin{pmatrix}
\alpha_s & 0
\end{pmatrix} = \alpha_s \otimes 
\begin{pmatrix}
1 & 0
\end{pmatrix}
=0
\end{equation*}
in $Q \otimes_{Q'} (\Omega_{Q'} W)_{\oplus}$, yet we still have
\begin{equation*}
\Lambda({\rm dot}_s)= 1 \otimes 
\begin{pmatrix}
\alpha_s & 0
\end{pmatrix} \neq 0
\end{equation*}
since $\begin{pmatrix} 1 & 0 \end{pmatrix} \notin \mathcal{M}'$. 
\end{rem}

\bibliographystyle{habbrv}
\bibliography{../quantummatrixrecursion-refs.bib}
\end{document}